\documentclass[11pt,letterpaper]{article}
\usepackage[T1]{fontenc}
\IfFileExists{lmodern.sty}{\usepackage{lmodern}\usepackage{microtype}}{\usepackage[expansion=false]{microtype}}\usepackage{amsmath,amssymb,amsthm,mathtools,booktabs,array,longtable}
\usepackage[margin=1in]{geometry}
\usepackage{xurl}
\usepackage[hidelinks]{hyperref}
\usepackage{fancyhdr}
\pagestyle{fancy}\fancyhf{}
\fancyhead[L]{Rational-circle obstructions for magic squares of squares}
\fancyfoot[C]{\thepage}
\setlength{\headheight}{14pt}
\setlength{\emergencystretch}{2em}
\setlength{\tabcolsep}{5pt}
\renewcommand{\arraystretch}{1.12}
\newtheorem{theorem}{Theorem}[section]
\newtheorem{proposition}[theorem]{Proposition}
\newtheorem{lemma}[theorem]{Lemma}
\newtheorem{corollary}[theorem]{Corollary}
\theoremstyle{definition}\newtheorem{definition}[theorem]{Definition}
\newtheorem{conjecture}[theorem]{Conjecture}
\theoremstyle{remark}\newtheorem{remark}[theorem]{Remark}
\newcommand{\Q}{\mathbb Q}\newcommand{\Z}{\mathbb Z}\newcommand{\R}{\mathbb R}
\newcommand{\F}{\mathbb F}\newcommand{\T}{\mathbb T}\newcommand{\PP}{\mathbf P^1}
\newcommand{\im}{\operatorname{Im}}\newcommand{\Res}{\operatorname{Res}}
\newcommand{\rk}{\operatorname{rank}}\newcommand{\diag}{\operatorname{diag}}
\newcommand{\Qsq}{\Q^{2}}
\hypersetup{pdftitle={Rational-circle obstructions and a genus-five reduction for magic squares of squares},pdfauthor={Kimia Samiein, Jonathan Wolinski}}
\title{Rational-circle rigidity for magic squares of squares:\\certified exclusions, a genus-five reduction,\\and surfaces of general type}
\author{Kimia Samiein \and Jonathan Wolinski}
\date{22 October 2026 (version 3: author-review draft with the 22 October additions)}
\begin{document}
\maketitle
\begin{abstract}
We study rational-circle systems arising from the four pairs of opposite entries of a $3\times3$ magic square of distinct positive integer squares. For a primitive central root $y=p^kq^\ell$, local valuations restrict the four frequency vectors to a two-parameter normal form. We give exact certificates excluding every such rational-circle configuration with $k+\ell\le7$. For actual integer squares the two-prime case is excluded outright by an arithmetic argument in the public Magic Squares of Squares archive, which we have verified; our theorem concerns the larger rational-circle problem, which that argument does not reach, and it implies the weaker arithmetic condition $k+\ell\ge8$ without Gaussian-prime arithmetic. The certificates include all signed assignments, exceptional fibers, and projective parameter values. For a three-parameter configuration, retaining the rational normalization of a unit-circle direction imposes two simultaneous square conditions. Sixteen of its forty-eight signed systems are immediately degenerate, and three elementary elliptic descents exclude another twenty-four. The remaining existence problem is equivalent to the existence of a rational point outside eight specified points on an explicit genus-five curve; rational reconstruction formulas prove the converse. Integer monomial maps transfer an exclusion of this configuration to three other configurations. The genus-five point list is not proved complete. We conjecture that no full-depth configuration in any number of parameters admits an admissible rational-circle point; this is strictly stronger than the nonexistence of a primitive magic square of distinct squares, and the results above are its first proven cases. For four parameters the solution sets are surfaces. An exhaustive search over all $11{,}862$ squarefree full-depth configurations finds no admissible point with three parameters of projective height at most four, and for the Hadamard configuration we prove that every one of its $48$ solution surfaces is of general type: each dominates, by a generically finite map of degree eight, a minimal surface with $p_g=11$, $q=0$, $K^2=32$, the bidouble cover of $\PP\times\PP$ branched along two bidegree-$(4,4)$ curves. The general-type behaviour known for the magic-square surface thus already appears in the trigonometric relaxation. The results concern specified subfamilies of that relaxation and do not decide the unrestricted magic-square problem. \emph{Additions of 22 October 2026.} The two-parameter certificate is extended through exponent sum nine ($5{,}520$ signed systems, independently re-checked). The genus-five curve has the involutions $h\mapsto-h$ and $h\mapsto1/h$; its three elliptic quotients are the curves 123.a1, 492.b1 and 1968.j1 of the LMFDB, each with a rational point of infinite order, so the eight-point conjecture is not reachable by rank-zero descent or by classical Chabauty on the genus-two quotient and is a quadratic-Chabauty problem. The torsion-coset hypothesis holds on all $662{,}528$ rational points of the solution varieties with three parameters of projective height at most four; a degeneracy-rank lemma explains the structure of these points and predicts in which configurations a point can lie on a torsion coset of codimension one. No rational point of height up to $10^6$ exists on the genus-five curve beyond the eight known.
\end{abstract}
\medskip
\noindent\textbf{Keywords.} Magic squares of squares; rational points; Gaussian integers; elliptic descent; exact computation.

\section{Introduction}\label{sec:introduction}
The magic-square-of-squares problem asks whether a $3\times3$ array of nine distinct positive integer squares can have equal sums along all three rows, all three columns, and the two long diagonals. Arithmetic formulations of this question are studied by Robertson~\cite{Robertson} and Bremner~\cite{Bremner}. Brown describes a formulation in terms of four rational right triangles on a common hypotenuse and proves a restriction arising from four sum-of-two-squares representations~\cite{Brown511,Brown417}. We use the common-hypotenuse formulation, keeping the distinction between integral triangles at a prescribed center and rational triangles up to rescaling.

Let $y^2$ be the central entry. Each opposite pair has one entry above and one below the center. The four positive differences between the larger entry of each pair and the center have the form
\begin{equation}\label{eq:pattern}
 \{a,b,a+b,a-b\},\qquad a>b>0,\quad a\ne2b.
\end{equation}
Each difference is $2pq$ for an integral right triangle with legs $p,q$ and hypotenuse $y$. Gaussian factorization then expresses the differences as sine coordinates of products of rational points on the unit circle. The use of Gaussian integers in this problem is developed by Cain~\cite{Cain} and Wolird~\cite{Wolird}; organization by the prime signature of the center is explicitly proposed by Weisenberg~\cite{Weisenberg}.

The circle formulation itself is also prior work. Houston~\cite[formulation P4 and question W1]{Houston} distinguishes the square subgroup of the rational circle from the full circle group and asks a weaker additive-coordinate question on the latter. Our objects are frequency-restricted subfamilies of that relaxation. The restrictions retain relations among four circle points while removing the requirement that their generators come from specified Gaussian primes.

\subsection{The rational-circle systems}
Write
\[
 \T(\Q)=\{z\in\Q(i):z\bar z=1\},
\]
with multiplication as its group law. Angle notation $z=e^{i\theta}$ will always mean that both $\cos\theta$ and $\sin\theta$ are rational; it does not mean that $\theta$ is a rational number. For an integer matrix $A$ with four rows $v_0,\ldots,v_3$, define
\[
 s_j(z_1,\ldots,z_r)=\im\prod_{a=1}^r z_a^{v_{ja}}.
\]
A point of $\T(\Q)^r$ is \emph{admissible for $A$} if the four magnitudes $|s_j|$ are distinct, nonzero, and have the pattern~\eqref{eq:pattern}. This condition is unchanged by permuting the rows or changing any row's sign.

The principal two-parameter configurations are
\begin{equation}\label{eq:normalformintro}
 A_{k,\ell;m,n}=
 \begin{pmatrix}k&\ell\\k&-\ell\\m&\ell\\k&n\end{pmatrix},
 \qquad k,\ell\in\Z_{>0},\quad m,n\in\Z,\quad |m|<k,\quad |n|<\ell.
\end{equation}
Their necessity for primitive integer squares is proved in Section~\ref{sec:arithmetic}.

\begin{theorem}[Certified two-parameter exclusion]\label{thm:main-two}
For $k+\ell\le7$, no matrix $A_{k,\ell;m,n}$ in~\eqref{eq:normalformintro} admits an admissible point of $\T(\Q)^2$. Consequently, if a primitive magic square has central root $y=p^kq^\ell$ with distinct primes $p,q$ and positive exponents, then $k+\ell\ge8$.
\end{theorem}
The arithmetic consequence is now known to be far from sharp: the public archive proves that no primitive square has a two-prime central root at all~\cite{MSSTwoBlock}, by an argument that uses the Gaussian-prime structure essentially, and we report our verification of that proof in Section~\ref{sec:prior}. The content of Theorem~\ref{thm:main-two} is the exclusion on the larger rational-circle domain, which that argument does not address; the corollary records what the relaxed statement alone yields. The conclusion is uniform in the rational-circle parameters, not a bounded search in the primes or in the center. Its proof consists of a finite exact elimination together with independently checkable certificates (Section~\ref{sec:two}). The same proof yields an obstruction for arbitrary exponents whenever the reduced denominators of $m/k$ and $n/\ell$ have sum at most seven. Theorem~\ref{thm:sumnine} (added 22 October 2026) extends the certificate through exponent sum nine, so ``seven'' may be read as ``nine'' in this paragraph.

Our second principal result concerns the three-parameter system
\begin{equation}\label{eq:canonicalintro}
 \sin(\alpha+\beta),\qquad\sin(\alpha-\beta),\qquad
 \sin\alpha,\qquad\sin\gamma.
\end{equation}
It is an invertible integer change of variables of the configuration labeled $D1$ in Section~\ref{sec:three}. Put
\begin{align}
 P(h)&=h^{12}+70h^{10}-1009h^8+7060h^6
                  -1009h^4+70h^2+1,\label{eq:Pintro}\\
 \mathcal C&:\quad Y^2=P(h).\notag
\end{align}
Here and below a rational-point statement about a hyperelliptic curve refers to its smooth projective model. Let
\begin{equation}\label{eq:knownintro}
 \mathcal K=\{(0,\pm1),(1,\pm72),(-1,\pm72),\infty_+,\infty_-\}.
\end{equation}

\begin{theorem}[Normalization and a genus-five reduction]\label{thm:main-d1}
The curve $\mathcal C$ has genus five. System~\eqref{eq:canonicalintro} has an admissible rational-circle solution if and only if $\mathcal C(\Q)\setminus\mathcal K$ is nonempty. The implication from a point outside $\mathcal K$ to a solution is given by explicit rational formulas.
\end{theorem}
Our third result concerns four parameters, where the relaxed solution sets are surfaces. We state a single conjecture, \emph{rational-circle rigidity} (Conjecture~\ref{conj:rigidity}), asserting that no full-depth configuration in any number of parameters admits an admissible rational-circle point; the theorems above are its first proven cases. An exhaustive search over all $11{,}862$ squarefree four-parameter configurations at projective height at most four finds nothing (Proposition~\ref{prop:four}), against the expectation that surfaces carry many rational points, and Theorem~\ref{thm:gentype} explains the discrepancy for the Hadamard configuration: all $48$ of its solution surfaces are of general type, with explicit invariants.

Theorem~\ref{thm:main-d1} is an existence equivalence, not a claimed bijection between the two parameter spaces. In particular, the theorem does not assert $\mathcal C(\Q)=\mathcal K$. The proof retains a rational-square normalization condition before eliminating variables. It excludes sixteen signed systems directly and twenty-four by three rank-zero elliptic descents, then treats the eight remaining systems through~\eqref{eq:Pintro} (Sections~\ref{sec:norm}--\ref{sec:genusfive}). Integer-torus maps show that $\mathcal C(\Q)=\mathcal K$ would also exclude three other configurations, labeled $B2,C,E2$.

\subsection{Relation to prior work and scope}\label{sec:prior}
The local valuation formula and the fact that at least three opposite-pair differences are units at each center prime appear in Sections~3--4 of the public \emph{Magic Squares of Squares} research archive~\cite{MSSArchive}. We cite that argument and include a proof to specify exactly where primitivity enters. A separate note in the archive claims a uniform two-prime exclusion for actual integer squares~\cite{MSSTwoBlock}, and a companion note claims a squarefree three-prime exclusion through a twenty-eight-class analysis. Neither claim is an assumption of this paper. We have independently verified the final section of the two-prime note — the coupled norm equations, their elimination to a single identity, the resulting quartic $y^2=t^4-10t^2+1$, and the rank-zero curve $y^2=x^3-x^2-2x$ of conductor $96$ with its two-isogeny descent — and the integrality mechanism on which the argument rests, namely that a rational integer divisible by $\pi^{4k}$ in $\Z[i]$ is divisible by $p^{4k}$; and we have re-derived the structural reductions on which it rests: in the block-balance note, the exact valuation spectrum, the three-unit lemma, the height bounds with their odd-quotient refinement, the uniqueness of the exception for a dominant block, the exclusion of balanced and midpoint indices, the exact corner-quotient valuation $v_p(\nu)=4h$ that rules out a dominant corner exception, and the modulo-$8$ refinement; and in the two-prime note, the reduction forcing both exceptions to be edges with both primes $\equiv1\pmod8$. We have also re-derived its balanced-exception case analysis (Sections 1--6) in full, including the four sign rows of its final mixed assignment and its quartic lemma, which is the reciprocal quartic of our Proposition~\ref{prop:quartics} reached through a $2$-isogenous model of conductor $56$. The only inputs we did not re-derive are Fermat's classical theorem that $X^2+Y^4=Z^4$ has no positive solution, and two facts imported from an earlier note of the archive that are special cases of our Lemma~\ref{lem:twoform} and Corollary~\ref{cor:parallelogram}: the two exceptions are distinct and cannot both be balanced. We therefore regard the archive's two-prime theorem as verified. Of the three-prime note (twenty-eight classes, with a transfer argument to arbitrary balanced or extreme exponents) we have verified every elliptic curve it relies on — $Y^2=X(X+1)(X+5)$ of conductor $80$, which closes its all-extreme class and is the curve of our orbit A; $Z^2=T(T+1)(T+9)$ of conductor $48$ with its complete point set $T\in\{0,-1,-9,\pm3\}$; and $Y^2=X^3\pm4X$-type curves of conductor $32$, all of rank zero — together with its first arithmetic class in full (which falls to Fermat's $X^2+Y^4=Z^4$ after a reduced-fraction argument), the block-generic facts of its transfer section, and the orientation-incoherence theorem it rests on in full: both the $g=1$ reduction (no primitive cube root of unity lies in $\Q(i)$) and the proofs that no coefficient-$111$ relation is coherent (four complex units of vanishing sum split into antipodal pairs, forcing two Gaussian blocks to be units) and no coefficient-$211$ relation is coherent (the reduction $2xy+x-y=2$ yields $R^2+S^2=A^2$ and $R^2+9S^2=W^2$ simultaneously, hence a positive rational square on $V^2=T(T+1)(T+9)$, of rank zero). Its per-class exclusions in Sections~6--16 and its exponent-transfer argument we have not re-derived line by line. Every claim of the archive we have checked has been correct. Consequently the arithmetic corollary of Theorem~\ref{thm:main-two} is subsumed by the archive's theorem; Theorem~\ref{thm:main-two} itself is not, since its domain is $\T(\Q)^2$ with specified frequencies, whereas the archive's argument uses the Gaussian-prime arithmetic at every step and says nothing about the relaxed problem.

The global geometry of the magic-square surface is studied by Bruin, Thomas, and V\'arilly-Alvarado~\cite{BTV}, and by Auel and Singer~\cite{AuelSinger}. The present paper does not compute invariants of that surface. Its contributions are the specified rational-circle exclusion theorem, explicit transfer maps between configurations, and the normalization-based equivalence in Theorem~\ref{thm:main-d1}. We make no priority claim for the common-hypotenuse dictionary, the rational-circle relaxation, the general monomial-map construction, or elliptic descent.

The exposition is confined to these results and their proof dependencies. Broader calculations on other three-parameter components and elliptic height bounds are retained in the accompanying extended research record, but are not needed for either main theorem. No completed Chabauty computation or classification of all three-parameter systems is claimed.

\section{The arithmetic source of the frequency restrictions}\label{sec:arithmetic}
We call an integer-square magic square \emph{primitive} when the greatest common divisor of its nine positive square roots is one. Dividing all entries by the square of that greatest common divisor preserves the magic conditions and distinctness.

\begin{lemma}[Lucas form]\label{lem:lucas}
Every rational $3\times3$ magic square has the form
\[
 \begin{pmatrix}
 e+a&e-a-b&e+b\\
 e-a+b&e&e+a-b\\
 e-b&e+a+b&e-a
 \end{pmatrix},\qquad\text{magic sum }3e.
\]
Its entries are pairwise distinct if and only if
\begin{equation}\label{eq:distinct}
 ab(a^2-b^2)(a^2-4b^2)(4a^2-b^2)\ne0.
\end{equation}
The magnitudes of its four opposite-pair differences can then be written as~\eqref{eq:pattern}.
\end{lemma}
\begin{proof}
Let the magic sum be $S$. Adding the four lines through the center counts every entry once and the center three additional times. Thus $4S=3S+3e$, giving $S=3e$. Every opposite pair sums to $2e$. Choosing the top-left and top-right differences $a,b$ determines all remaining entries from the row and column equations. The nine differences are $0,\pm a,\pm b,\pm(a+b),\pm(a-b)$. Their pairwise distinctness is exactly~\eqref{eq:distinct}. Replacing $a,b$ by their ordered absolute values gives~\eqref{eq:pattern}.
\end{proof}

\begin{proposition}[Integral common-hypotenuse dictionary]\label{prop:dictionary}
For a fixed positive integer $y$, a magic square of nine distinct positive integer squares with center $y^2$ exists if and only if there are four integral right triangles $(p_j,q_j,y)$ whose positive numbers $2p_jq_j$ form~\eqref{eq:pattern}.
\end{proposition}
\begin{proof}
If $r^2,y^2,s^2$ form an increasing progression of integer squares, then $r^2+s^2=2y^2$. The roots $r,s$ have the same parity, and
\[
 p=\frac{s+r}{2},\qquad q=\frac{s-r}{2}
\]
are positive integers satisfying $p^2+q^2=y^2$ and $s^2-y^2=2pq$. Conversely, an integral right triangle gives the progression
\[
 (p-q)^2,\qquad y^2,\qquad(p+q)^2.
\]
Its first term is positive: nonzero equal integer legs would imply that an integer square is twice another nonzero integer square. Apply the correspondence to all four center lines and use Lemma~\ref{lem:lucas}. In the converse, the largest difference $a+b$ is less than $y^2$, since it is one of the triangle differences, so all nine reconstructed entries are positive.
\end{proof}
We call $2pq$ the \emph{offset}. It is four times the ordinary triangle area. Brown's common-hypotenuse formulation~\cite{Brown511} permits rational triangles when scale is allowed to vary. The integral formulation above is required when the integer $y$, and hence its factorization, is held fixed.

\begin{lemma}[Prime support]\label{lem:support}
In a primitive square, $y$ is odd and every prime divisor of $y$ is congruent to $1$ modulo $4$.
\end{lemma}
\begin{proof}
For each opposite pair, $r^2+s^2=2y^2$. If $2\mid y$, reduction modulo $4$ forces $r,s$ even. If an odd prime $p\equiv3\pmod4$ divides $y$, reduction modulo $p$ forces both roots divisible by $p$, since $-1$ is not a square modulo $p$. In either case the prime divides all nine roots, contrary to primitivity.
\end{proof}

Write
\[
 y=\prod_{j=1}^r\ell_j^{k_j},\qquad
 \ell_j=\pi_j\overline{\pi_j}\quad\text{in }\Z[i],
\]
where the $\ell_j$ are distinct primes congruent to $1$ modulo $4$, the $k_j$ are positive integers, and one Gaussian prime above each $\ell_j$ is fixed. Unique factorization gives the following indexing.
\begin{proposition}[Gaussian coordinates]\label{prop:gaussian}
The nondegenerate integral right triangles on $y$, up to interchanging legs, are indexed by
\[
 0\ne m=(m_1,\ldots,m_r)\in\Z^r,\qquad |m_j|\le k_j,
 \qquad m\sim-m.
\]
A representative is
\begin{equation}\label{eq:gaussian}
 Z_m=\prod_j\pi_j^{k_j+m_j}\overline{\pi_j}^{k_j-m_j},
 \qquad Z_m\overline{Z_m}=y^2.
\end{equation}
The legs are $|\operatorname{Re}Z_m|$ and $|\operatorname{Im}Z_m|$, in either order. Its primitive hypotenuse and common leg divisor are
\[
 d_m=\prod_j\ell_j^{|m_j|},\qquad
 g_m=\prod_j\ell_j^{k_j-|m_j|}.
\]
There are $\bigl(\prod_j(2k_j+1)-1\bigr)/2$ such triangles.
\end{proposition}
\begin{proof}
The exponents of $\pi_j$ and $\bar\pi_j$ in a Gaussian integer of norm $y^2$ add to $2k_j$, giving~\eqref{eq:gaussian} up to a unit. Units and conjugation account for signs of the legs and leg interchange; conjugation sends $m$ to $-m$. An axis representation has equal exponents above every $\ell_j$, so only $m=0$ is excluded. Equal nonzero legs are impossible over the integers. Removing the common rational factor $g_m$ leaves only one of $\pi_j,\bar\pi_j$ above each prime in the primitive Gaussian representative, so there is no remaining common leg divisor. The count follows.
\end{proof}

The angle encoding an offset is
\begin{equation}\label{eq:areaangle}
 e^{i\chi_j}=\left(\frac{\pi_j}{\bar\pi_j}\right)^2
             =\frac{\pi_j^4}{\ell_j^2}\in\T(\Q).
\end{equation}
Indeed, $Z_m/y=\prod_j(\pi_j/\bar\pi_j)^{m_j}$, so
\begin{equation}\label{eq:offset}
 2p_mq_m=|\im Z_m^2|
       =y^2\left|\sin\left(\sum_jm_j\chi_j\right)\right|.
\end{equation}
The square in~\eqref{eq:areaangle} is essential. Replacing the particular points there by arbitrary elements of $\T(\Q)$ enlarges the problem. For instance, $i\in\T(\Q)$ has no square root in $\T(\Q)$. This is the distinction between the circle and its square subgroup emphasized by Houston~\cite{Houston}.

\begin{lemma}[Valuations and full depth]\label{lem:depth}
Let $(p,q,y)=g(p',q',d)$ be an integral right triangle with $(p',q',d)$ primitive and $gd=y$. If $\ell^k\Vert y$, then
\[
 v_\ell(2pq)=2\bigl(k-v_\ell(d)\bigr)\quad(\ell\mid d),
 \qquad v_\ell(2pq)\ge2k\quad(\ell\nmid d).
\]
In a primitive magic square, for each prime power $\ell^k\Vert y$, at least three of the four triangle vectors satisfy $|m_\ell|=k$. The three vectors may depend on $\ell$.
\end{lemma}
\begin{proof}
The hypotenuse of a primitive integer right triangle is odd, and $\gcd(2p'q',d)=1$. The valuation formulas follow from $2pq=g^2(2p'q')$.

For odd $\ell$, the minimum valuation among $a,b,a+b,a-b$ occurs at least three times. Unequal valuations of $a,b$ make both sums have the smaller valuation. If the two valuations are equal, both sums cannot have a larger valuation, since their sum is $2a$. In a primitive square the common minimum is zero: otherwise $\ell\mid y^2$ and every offset would make every square entry, hence every root, divisible by $\ell$. The formula and Proposition~\ref{prop:gaussian} identify valuation zero with $v_\ell(d)=|m_\ell|=k$.
\end{proof}
The valuation and three-unit argument are recorded in~\cite[Sections 3--4]{MSSArchive}. The proof above is included to make their integral and primitive hypotheses explicit.

\begin{corollary}\label{cor:primepower}
The central root of a primitive integer-square magic square is not a prime power.
\end{corollary}
\begin{proof}
For $y=\ell^k$, there is only one projective vector of full depth, represented by $m=k$. Lemma~\ref{lem:depth} requires three distinct full-depth triangles.
\end{proof}

\begin{lemma}[Two-prime normal form]\label{lem:twoform}
If $y=p^kq^\ell$ is the central root of a primitive square, the four triangle vectors, modulo individual signs, have the form~\eqref{eq:normalformintro}.
\end{lemma}
\begin{proof}
At least three vectors have first coordinate of magnitude $k$, and at least three have second coordinate of magnitude $\ell$. Hence at least two have both magnitudes maximal. Only two projective vectors have this property, namely $(k,\ell)$ and $(k,-\ell)$, so both occur. A third cannot occur. The deficient first and second coordinates must therefore be on different remaining vectors. Change each vector's sign to make its full coordinate positive.
\end{proof}
These restrictions are derived for actual primitive squares before the rational-circle parameters are relaxed. They are not consequences of the unrestricted circle equations alone.

\section{Signed equations and an elementary obstruction}\label{sec:equations}
\begin{lemma}[Complete signed assignments]\label{lem:roles}
For real $s_0,s_1,s_2,s_3$ with nonzero pairwise-distinct magnitudes, admissibility is equivalent to at least one system
\begin{equation}\label{eq:roles}
 s_i-\varepsilon s_k-\eta s_l=0,\qquad
 s_j-\varepsilon s_k+\eta s_l=0,
\end{equation}
where $i<j$, $(k,l)$ is either order of the complementary pair, and $\varepsilon,\eta\in\{1,-1\}$. There are $6\cdot2\cdot4=48$ such systems.
\end{lemma}
\begin{proof}
Choose $i,j$ for the magnitudes $a+b,a-b$. Whatever their signs and order, the magnitudes of $(s_i+s_j)/2$ and $(s_i-s_j)/2$ are $a,b$ in some order. Choose $k,l$ in that order and select $\varepsilon,\eta$. Conversely,~\eqref{eq:roles} gives
\[
 \{|s_i|,|s_j|\}
 =\{|s_k|+|s_l|,\,\bigl||s_k|-|s_l|\bigr|\}.
\]
Take $a=\max(|s_k|,|s_l|)$ and $b=\min(|s_k|,|s_l|)$. Distinctness and nonvanishing give the restrictions in~\eqref{eq:pattern}.
\end{proof}

\begin{theorem}[A real-circle obstruction]\label{thm:real}
For arbitrary real $x,z$, the four magnitudes
\[
 X=|\sin x|,\quad Z=|\sin z|,\quad
 P=|\sin(x+z)|,\quad Q=|\sin(x-z)|
\]
are not admissible.
\end{theorem}
\begin{proof}
Addition and subtraction of sines give
\begin{align}
 P^2Q^2&=(X^2-Z^2)^2,\label{eq:product}\\
 P^2+Q^2&=2(X^2+Z^2-2X^2Z^2).\label{eq:sum}
\end{align}
Assign the four values in~\eqref{eq:pattern} to $X,Z,P,Q$. There are six choices for the unordered pair $\{P,Q\}$. Substitution in~\eqref{eq:product} gives the following residuals:
\begin{center}
\begin{tabular}{ll}
\toprule
$\{P,Q\}$ & $P^2Q^2-(X^2-Z^2)^2$\\
\midrule
$\{a,b\}$ & $-15a^2b^2$\\
$\{a,a+b\}$ & $3a^2b(2a-b)$\\
$\{a,a-b\}$ & $-3a^2b(2a+b)$\\
$\{b,a+b\}$ & $-3ab^2(a-2b)$\\
$\{b,a-b\}$ & $-3ab^2(a+2b)$\\
$\{a+b,a-b\}$ & $0$\\
\bottomrule
\end{tabular}
\end{center}
The first five are nonzero under~\eqref{eq:pattern}. In the last case $\{X,Z\}=\{a,b\}$, and the residual in~\eqref{eq:sum} is $4a^2b^2\ne0$.
\end{proof}

\begin{corollary}\label{cor:parallelogram}
A configuration $\{v,w,v+w,v-w\}$, up to individual signs, has no admissible point even over the real circle. In particular, no primitive square has central root $pq$ for distinct primes $p,q$.
\end{corollary}
\begin{proof}
Substitute $x=v\cdot\theta$, $z=w\cdot\theta$ in Theorem~\ref{thm:real}. For $(k,\ell)=(1,1)$, Lemma~\ref{lem:twoform} gives exactly this configuration.
\end{proof}
The semiprime arithmetic conclusion is contained in Brown's Proposition~1~\cite{Brown417}. The proof above is stated over the real circle because that form will also be used in the parameter-space reductions.

Every element of $\T(\Q)$ can be written uniquely in the projective parametrization
\begin{equation}\label{eq:circleparam}
 z=\frac{(b+ia)^2}{b^2+a^2},\qquad[a:b]\in\PP(\Q).
\end{equation}
The inverse is $a/b=\im z/(1+\operatorname{Re}z)$ when $z\ne-1$; the point $z=-1$ corresponds to $[1:0]$. For a frequency $v=(v_1,\ldots,v_r)\in\Z^r$ with $|v_j|\le k_j$, a common homogeneous numerator of its sine is
\begin{equation}\label{eq:homogeneous}
 E_v=\im\prod_j
 (b_j+i\,\operatorname{sgn}(v_j)a_j)^{2|v_j|}
 (b_j^2+a_j^2)^{k_j-|v_j|}.
\end{equation}
The denominator $\prod_j(b_j^2+a_j^2)^{k_j}$ is positive at every real projective point. The affine parameter $a_j/b_j$ and the point $[1:0]$ are treated separately below.

\section{Certified exclusion through exponent sum seven}\label{sec:two}
\subsection{Reduction of the configuration list}
\begin{lemma}[Integer-power inheritance]\label{lem:inherit}
Let a four-row integer frequency matrix $A$ have no admissible rational-circle point. Multiplying any column by a positive integer preserves this exclusion. Coordinate permutations, coordinate sign changes, and signed row permutations also preserve it.
\end{lemma}
\begin{proof}
If the column multipliers are $d_j$, a solution for the multiplied matrix at $(z_j)$ gives a solution for $A$ at $(z_j^{d_j})$. The latter points lie in $\T(\Q)$ and have the same four sine magnitudes. The other transformations either relabel the parameters or change individual sine signs.
\end{proof}
Only the forward power map is used. A rational point need not have a rational preimage under a nontrivial power map.

After interchanging the two coordinates, assume $k\ge\ell$. There are twelve positive exponent pairs with $k+\ell\le7$, containing
\[
 \sum_{k\ge\ell>0,\ k+\ell\le7}(2k-1)(2\ell-1)=168
\]
raw normal forms. The cases $m=n=0$ follow from Corollary~\ref{cor:parallelogram}. Lemma~\ref{lem:inherit} and the stated symmetries reduce the rest to the forty-three representatives in Appendix~\ref{app:representatives}. The checker enumerates every raw form and records its representative and column multipliers; the totals alone are not used as evidence of coverage.

\subsection{Polynomial inputs and proof certificates}
Define $C_j(v),S_j(v)$ for $j\ge0$ by
\[
 C_j(v)+iS_j(v)=(1+iv)^{2j},
 \qquad C_{-j}=C_j,\quad S_{-j}=-S_j.
\]
Put $t=\tan(x/2)$ and $u=\tan(z/2)$. On the denominator
\[
 L=(1+t^2)^k(1+u^2)^\ell,
\]
the four signed numerators for $A_{k,\ell;m,n}$ are
\begin{align}
 E_0&=S_k(t)C_\ell(u)+C_k(t)S_\ell(u),\notag\\
 E_1&=S_k(t)C_\ell(u)-C_k(t)S_\ell(u),\notag\\
 E_2&=(1+t^2)^{k-|m|}\bigl(S_m(t)C_\ell(u)+C_m(t)S_\ell(u)\bigr),\label{eq:twoinputs}\\
 E_3&=(1+u^2)^{\ell-|n|}\bigl(S_k(t)C_n(u)+C_k(t)S_n(u)\bigr).\notag
\end{align}
The first two are the numerators of $\sin(kx+\ell z)$ and $\sin(kx-\ell z)$; the other two correspond to the deficient vectors. The denominator is positive in the finite real chart.

\begin{lemma}[Exact elimination certificate]\label{lem:elimination}
For each of the forty-three representatives and each of the forty-eight signed systems in Lemma~\ref{lem:roles}, remove only a common power of $t$ from the two polynomial equations, when present. The resultant in $u$ of the resulting equations is nonzero. Its finite rational zeros are contained in $\{0,1,-1\}$.
\end{lemma}
\begin{proof}[Computer-assisted proof]
The certificate \path{verification/two_prime/results/certificate_degree7.json} contains the 2,064 equation pairs, every removed power, the reduced equations, the exact resultant, and its complete factor product. The factor catalog contains three linear polynomials with roots $0,1,-1$ and 1,187 nonlinear polynomials. Each nonlinear factor has a witness prime in
\[
 \{3,5,7,11,13,17,19\}
\]
at which its leading coefficient is nonzero and it has no root. Such a polynomial has no rational root: a rational root whose reduced denominator is divisible by the witness prime contradicts the leading coefficient after clearing denominators; a denominator coprime to the prime gives a root in the finite field.

All factor products, constants, and multiplicities are checked exactly. No irreducibility claim is required. For example, the factor $t^4-22t^2+25$ reduces modulo $3$ to $t^4+2t^2+1$, which takes the value $1$ at each residue. Its leading coefficient is nonzero, so the factor has no rational zero.

The first implementation reconstructs~\eqref{eq:twoinputs} by repeated sine addition. For the reduced equations $F,G$, put
\[
 M=\deg_uF,\quad N=\deg_uG,\quad A=\deg_tF,\quad B=\deg_tG.
\]
The Sylvester determinant has degree at most $D=NA+MB$ in $t$, since it uses $N$ coefficient rows from $F$ and $M$ from $G$. The checker rejects a recorded resultant of degree exceeding $D$ and compares it with the exact integer determinant at $t=0,1,\ldots,D$. These $D+1$ comparisons prove the polynomial identity. Fraction-free elimination is used with checked exact divisions. Across the certificate, 132,336 integer determinants are evaluated.

The second implementation reconstructs the numerators with Chebyshev polynomials and evaluates the symbolic Sylvester determinants over $\Z[t]$. Neither checker imports the certificate producer; neither uses factorization or a rational-root solver to accept the factor catalog. Both checks succeed for all 2,064 systems. Their exact inputs and outputs are part of the accompanying package. Appendix~\ref{app:representatives} gives a complete worked system, its Sylvester determinant, and a modular no-root witness.

Every finite common zero of $F,G$ annihilates the resultant, including points where a leading coefficient specializes to zero. The factor witnesses therefore prove the stated list of possible rational coordinates. The removed $t=0$ branches are handled in the next paragraph.
\end{proof}

\begin{proof}[Proof of Theorem~\ref{thm:main-two}]
Coverage reduces every normal form in the stated range to the real-circle obstruction or to a column-power pullback of a representative in Lemma~\ref{lem:elimination}. For a representative, every finite rational common zero has $t=0,1,-1$. The first two signed values satisfy
\begin{equation}\label{eq:collision}
 \sin^2(kx+\ell z)-\sin^2(kx-\ell z)
 =\sin(2kx)\sin(2\ell z).
\end{equation}
At these values of $t$, one has $x=0,\pm\pi/2$ modulo $2\pi$, and the right side vanishes. Thus two magnitudes coincide. This also excludes every removed $t=0$ branch.

At the projective boundary $t=\infty$, $x=\pi$, and at $u=\infty$, $z=\pi$;~\eqref{eq:collision} excludes both boundaries and their intersection. Hence there is no admissible point anywhere on $\PP(\Q)^2$. Lemma~\ref{lem:inherit} transfers the exclusions, and Lemma~\ref{lem:twoform} gives the arithmetic consequence.
\end{proof}

\begin{corollary}[Reduced-denominator restriction]\label{cor:denominator}
An admissible rational-circle point for $A_{k,\ell;m,n}$ requires
\[
 \frac{k}{\gcd(k,|m|)}+\frac{\ell}{\gcd(\ell,|n|)}\ge8,
\]
where $\gcd(k,0)=k$.
\end{corollary}
\begin{proof}
Divide the columns by $d=\gcd(k,|m|)$ and $e=\gcd(\ell,|n|)$. The reduced normal form has bounds $K=k/d$, $L=\ell/e$. If $K+L\le7$, Theorem~\ref{thm:main-two} excludes it, and Lemma~\ref{lem:inherit} excludes the original form.
\end{proof}
The corollary concerns individual frequency configurations. It does not exclude every two-prime central root with larger exponent sum.

\subsection{Extension through exponent sum nine (22 October 2026)}\label{sec:sumnine}
The producer and checker of Lemma~\ref{lem:elimination} were rerun, unchanged in method, on the signatures with $k+\ell=8$ and $k+\ell=9$. Coverage was planned exactly as in Section~\ref{sec:two}: of the $140$ raw normal forms with $k+\ell=8$, $72$ are inherited from the sum-seven certificate by Lemma~\ref{lem:inherit} and the remaining $68$ form $32$ orbits; of the $172$ raw forms with $k+\ell=9$, $92$ are inherited and the remaining $80$ form $40$ orbits. The orbit decomposition is asserted exactly by the planning script and re-derived independently by the checker.

\begin{theorem}[Certified exclusion through exponent sum nine]\label{thm:sumnine}
For $k+\ell\le9$, no matrix $A_{k,\ell;m,n}$ in~\eqref{eq:normalformintro} admits an admissible point of $\T(\Q)^2$. Consequently Corollary~\ref{cor:denominator} holds with $8$ replaced by $10$, and a primitive magic square with central root $p^kq^\ell$ requires $k+\ell\ge10$.
\end{theorem}
\begin{proof}[Computer-assisted proof]
The merged certificate contains $115$ representatives and $5{,}520$ signed systems ($2{,}064$ from Lemma~\ref{lem:elimination}, $1{,}536$ new at sum eight, $1{,}920$ new at sum nine), a factor catalog of $3{,}207$ polynomials of which $3{,}204$ carry a modular no-root witness and three are the linear factors with roots $0,\pm1$, and $553{,}392$ exact integer Sylvester determinants. The largest witness prime is still $19$. Every finite rational fiber has $t\in\{0,\pm1\}$, and the projective boundaries are excluded by~\eqref{eq:collision} as before. The standard-library checker accepts the merged certificate in strict mode; the sum-eight part was additionally accepted by the symbolic Chebyshev/Sylvester checker and by the mutation tests, while the symbolic cross-check and mutation tests of the sum-nine part had not been run at the time of writing. Certificates, checkers, logs and the agent transcripts that produced them are in the companion repository (Section~\ref{sec:repro}).
\end{proof}
The arithmetic consequence $k+\ell\ge10$ is, like $k+\ell\ge8$, subsumed by the archive's two-prime theorem (Section~\ref{sec:prior}); the content of Theorem~\ref{thm:sumnine} is the relaxed-problem exclusion, which that argument does not reach.

\section{Three-parameter configurations and transport}\label{sec:three}
For a squarefree central root $y=\ell_1\ell_2\ell_3$, Proposition~\ref{prop:gaussian} gives thirteen nonzero projective vectors in $\{-1,0,1\}^3$. The full-depth condition requires at least three of the four selected vectors to have nonzero entry in each coordinate.

\begin{proposition}[Finite configuration list]\label{prop:orbits}
There are 177 four-element sets with this property. Under coordinate sign changes and coordinate permutations, they form fifteen orbits.
\end{proposition}
\begin{proof}[Finite enumeration]
Choose the sign of each nonzero vector by requiring its first nonzero coordinate to be positive. Enumerate the four-subsets of the resulting thirteen vectors and retain those satisfying the coordinate condition. Form their orbits under the $2^3\cdot3!=48$ signed coordinate permutations. The orbit representatives and sizes are listed in Appendix~\ref{app:orbits}; their sizes sum to 177. The independent enumeration in \path{check_paper.py} computes the vector set, the retained subsets, and every orbit, and verifies that the representatives exhaust the list.
\end{proof}
This is a classification under the stated finite group, not under all integer changes of circle coordinates. A larger group can identify configurations in different rows of the table.

For the transport statements, use these ordered frequency matrices:
\begin{equation}\label{eq:selectedmatrices}
\begin{aligned}
 A_{D1}&=\begin{pmatrix}1&1&1\\1&1&-1\\1&1&0\\1&0&1\end{pmatrix},&
 A_{B2}&=\begin{pmatrix}1&1&1\\1&1&-1\\1&-1&-1\\1&1&0\end{pmatrix},\\[2pt]
 A_C&=\begin{pmatrix}1&1&1\\1&1&-1\\1&-1&1\\0&0&1\end{pmatrix},&
 A_{E2}&=\begin{pmatrix}1&-1&1\\1&-1&-1\\1&1&0\\0&0&1\end{pmatrix}.
\end{aligned}
\end{equation}
Labels are identifiers for these explicit matrices, not additional hypotheses.

\begin{lemma}[Integer-torus transport]\label{lem:transport}
Suppose $A_sM=PA_t$, where $M$ is an integer square matrix of nonzero determinant and $P$ is a signed permutation matrix. An admissible point for $A_t$ maps to one for $A_s$ under the monomial map represented in angle notation by $\psi=M\theta$. If $\det M=\pm1$, this map has an integer inverse and gives an equivalence.
\end{lemma}
\begin{proof}
Integer products and inverse products preserve $\T(\Q)$. At the resulting point,
\[
 \sin(A_s\psi)=\sin(A_sM\theta)=P\sin(A_t\theta).
\]
The magnitudes are preserved up to permutation. When $M$ is unimodular its inverse is integral and supplies the converse.
\end{proof}

The following matrices give explicit transfers. The notation $M_{s,t}$ means that a point of configuration $t$ is mapped to one of configuration $s$:
\begin{align}
 M_{B2,C}&=\begin{pmatrix}1&0&1\\-1&0&0\\1&1&0\end{pmatrix},\notag\\
 M_{D1,B2}&=\begin{pmatrix}1&-1&-2\\0&2&2\\0&0&1\end{pmatrix},\label{eq:maps}\\
 M_{D1,C}&=\begin{pmatrix}0&-2&1\\0&2&0\\1&1&0\end{pmatrix},\notag\\
 M_{D1,E2}&=\begin{pmatrix}0&2&0\\0&-2&1\\1&-1&0\end{pmatrix}.\notag
\end{align}
Their frequency identities, including the signs and orderings of the output sine values, appear in Table~\ref{tab:transport}.
\begin{table}[htbp]
\centering
\begin{tabular}{ccrl}
\toprule
Source $s$ & Target $t$ & $\det M_{s,t}$ & $\sin(A_sM_{s,t}\theta)$ in target ordering\\
\midrule
$B2$ & $C$ & $-1$ & $(s_0,-s_1,s_2,s_3)$\\
$D1$ & $B2$ & $2$ & $(s_0,s_1,s_3,s_2)$\\
$D1$ & $C$ & $-2$ & $(s_0,-s_1,s_3,s_2)$\\
$D1$ & $E2$ & $2$ & $(s_0,-s_1,s_3,s_2)$\\
\bottomrule
\end{tabular}
\caption{Exact transfer identities. A target point maps to a source point.}\label{tab:transport}
\end{table}


Every identity in that table is a direct integer matrix product. The first matrix is unimodular; the others have determinant of absolute value two. Thus $B2$ and $C$ are equivalent rational-circle problems. A nonexistence theorem for $D1$ excludes $B2,C,E2$, but an arbitrary $D1$ point is not asserted to lift rationally through a determinant-two map.

\section{Rational normalization in the D1 system}\label{sec:norm}
For $A_{D1}$ in~\eqref{eq:selectedmatrices}, make the unimodular change
\[
 \alpha=x+z,\qquad\beta=w,\qquad\gamma=x+w.
\]
Its inverse is $x=\gamma-\beta$, $z=\alpha-\gamma+\beta$, $w=\beta$. The signed values become~\eqref{eq:canonicalintro}.

Put $r=\tan(\beta/2)$ and $q=1+r^2$. At a finite parameter,
\begin{equation}\label{eq:Hmatrix}
 \begin{pmatrix}s_0\\s_1\\s_2\\s_3\end{pmatrix}
 =q^{-1}H(r)\begin{pmatrix}\sin\alpha\\\cos\alpha\\\sin\gamma\end{pmatrix},
 \qquad
 H(r)=\begin{pmatrix}
 1-r^2&2r&0\\1-r^2&-2r&0\\q&0&0\\0&0&q
 \end{pmatrix}.
\end{equation}
The fibers $r=0,\infty$ are degenerate because the first two magnitudes coincide. Using sine and cosine coordinates for $\alpha,\gamma$ retains all of their rational-circle points, including the points omitted by their affine half-angle charts.

For each signed assignment, let $R$ be its two-by-four relation matrix from~\eqref{eq:roles}, and put $B(r)=RH(r)$. Write the cross product of its two rows as
\[
 g(r)v(r),\qquad v=(v_1,v_2,v_3),
\]
where $g\in\Z[r]\setminus\{0\}$ and $v\in\Z[r]^3$ have the normalization fixed by the exact record. The entries of $v$ have gcd one over $\Q[r]$. They therefore do not vanish simultaneously at any parameter. The only nonconstant factors of $g$ are
\[
 r,\qquad r^2-3,\qquad 3r^2-1.
\]
The last two have no rational root. Thus $B(r)$ has rank two at every nonzero rational $r$.

Every solution has
\[
 (\sin\alpha,\cos\alpha,\sin\gamma)=\lambda(v_1,v_2,v_3),
 \qquad \lambda\in\Q^\times.
\]
The two unit-circle equations imply
\begin{equation}\label{eq:twonorms}
 v_1^2+v_2^2\in\Qsq,
 \qquad
 v_1^2+v_2^2-v_3^2\in\Qsq.
\end{equation}
The first condition is the rational normalization of the direction $(v_1:v_2)$; the second is needed for a rational value of $\cos\gamma$. Neither may be dropped. Here $\Qsq=\{z^2:z\in\Q\}$ includes zero; the first quantity is nonzero at any solution.

\begin{lemma}[The normalization catalog]\label{lem:normcatalog}
For the forty-eight signed assignments, the identity
\[
 v_1^2+v_2^2=q^2N(r)
\]
has the complete list in Table~\ref{tab:normcatalog}.
\end{lemma}
\begin{table}[htbp]
\centering
\begin{tabular}{lrl}
\toprule
$N(r)$ & Assignments & Reduction\\
\midrule
$1$ &16& $v_1=0$\\
$r^4+6r^2+1$ &4& quartic\\
$r^4-3r^2+4$ &4& quartic\\
$4r^4-3r^2+1$ &4& reciprocal quartic\\
$r^4+18r^2+1$ &12& quartic\\
$r^2+9$ &4& conic and sextic\\
$9r^2+1$ &4& reciprocal conic and sextic\\
\bottomrule
\end{tabular}
\caption{Rational normalization conditions for all signed $D1$ assignments.}\label{tab:normcatalog}
\end{table}
\begin{proof}[Exact finite calculation]
For each relation matrix $R$, form $RH$ and its row cross product. Divide by the recorded common factor and verify the three resulting entries and their norm. The forty-eight polynomial vectors, the divisors $g$, the two norms in~\eqref{eq:twonorms}, and the assignment tuples are recorded in \path{verification/d1_norm/d1_certificate.json}. The checker reconstructs the matrices and polynomial products over $\Z[r]$, and verifies a constant gcd of the entries of $v$. The possible nonzero roots of $g$ are excluded as above. A second implementation computes the matrices symbolically and checks the historical frequency identities. Both produce exactly the counts in the table.

For example, the assignment $s_0=s_1+s_3$, $s_2=s_1-s_3$ gives $g=2r$ and
\[
 v=(3q,-rq,-4r^2).
\]
Its two norms are $q^2(r^2+9)$ and $r^6-5r^4+19r^2+9$. This is one of the residual conic cases.
\end{proof}

\subsection{Three quartic obstructions}
We use the classical descent associated with a rational point of order two; see Silverman and Tate~\cite{SilvermanTate} and the explicit 2-isogeny reference~\cite{SageDescent}. For
\[
 E:y^2=x^3+Ax^2+Bx,\qquad B(A^2-4B)\ne0,
\]
the 2-isogenous curve is
\[
 E':y^2=x^3-2Ax^2+(A^2-4B)x.
\]
The descent maps to $\Q^\times/\Q^{\times2}$ satisfy
\begin{equation}\label{eq:descentrank}
 2^{\rk E(\Q)+2}=|\alpha(E)|\,|\alpha'(E')|.
\end{equation}
For a point with $x\ne0$, $\alpha$ is the square class of $x$. An image class is represented by a signed squarefree divisor $d$ of $B$ and must permit an integral solution, with $\gcd(U,V)=1$, of
\begin{equation}\label{eq:cover}
 N^2=dU^4+A U^2V^2+(B/d)V^4.
\end{equation}
The identity and the rational two-torsion point are included in these classes with the usual definitions $\alpha(\mathcal O)=1$, $\alpha((0,0))=B$ modulo squares.

\begin{proposition}[Complete quartic point lists]\label{prop:quartics}
Each of
\[
 Z^2=r^4+6r^2+1,\quad
 Z^2=r^4-3r^2+4,\quad
 Z^2=r^4+18r^2+1
\]
has only its two points with $r=0$ and its two points at infinity. The reciprocal curve $Z^2=4r^4-3r^2+1$ likewise has no finite rational point with $r\ne0$.
\end{proposition}
\begin{proof}
The map $(r,Z)\mapsto(X,Y)=(r^2,rZ)$ from $Z^2=r^4+Ar^2+B$ lands on $E_{A,B}:Y^2=X^3+AX^2+BX$. We prove the complete groups in Table~\ref{tab:groups}.
\begin{table}[htbp]
\centering\small
\begin{tabular}{ccl}
\toprule
$(A,B)$ & Partner $(A',B')$ & $E_{A,B}(\Q)$\\
\midrule
$(6,1)$ & $(-12,32)$ & $\{\mathcal O,(0,0),(-1,2),(-1,-2)\}$\\
$(-3,4)$ & $(6,-7)$ & $\{\mathcal O,(0,0),(2,2),(2,-2)\}$\\
$(18,1)$ & $(-36,320)$ & $\{\mathcal O,(0,0),(-1,4),(-1,-4)\}$\\
\bottomrule
\end{tabular}
\caption{The three elliptic groups used in the normalization argument.}\label{tab:groups}
\end{table}
For $(6,1)$ and $(18,1)$ the first descent image is contained in $\{1,-1\}$. For $(-3,4)$ its negative classes are excluded over $\R$, leaving $\{1,2\}$. On the three partner curves the possible surviving classes are, respectively,
\[
 \{1,2\},\qquad\{-7,1\},\qquad\{1,5\}.
\]
For the first partner, negative classes make the right side of~\eqref{eq:cover} negative. For the second, classes $-1$ and $7$ are impossible modulo $16$. For the third, negative classes are impossible over $\R$, while $2$ and $10$ are impossible modulo $64$. The complete residue checks are specified in Appendix~\ref{app:descent}. Hence both image sizes are at most two in all three cases, and~\eqref{eq:descentrank} forces rank zero.

Good-prime point counts have greatest common divisor four: for the three models one may use $(p,\#E(\F_p))=(3,4),(5,8)$; $(3,4),(5,4)$; and $(3,4),(7,4)$, respectively. Prime-to-$p$ injectivity of torsion reduction bounds each primary component; for these particular counts it gives torsion order dividing four. Appendix~\ref{app:descent} supplies the prime-by-prime argument. The four displayed points therefore exhaust each group.

A finite nonzero rational $r$ would give a nonzero rational-square coordinate $X=r^2$. The only available nonzero $X$ is $-1$ in the first and third groups, or $2$ in the second group; neither is a nonzero rational square. This excludes all such $r$. Substitution gives $(0,\pm1)$, $(0,\pm2)$, and $(0,\pm1)$, respectively, and the square leading coefficient gives two rational infinity points. The polynomials are squarefree. Finally $r\mapsto1/r$ and $Z\mapsto Z/r^2$ identify the reciprocal quartics, including their projective points.
\end{proof}

\begin{corollary}\label{cor:forty}
Forty of the forty-eight signed $D1$ systems have no admissible rational-circle point. Every remaining solution, after possibly replacing $r$ by $1/r$, satisfies
\begin{equation}\label{eq:residual}
 S^2=r^2+9,\qquad
 Z^2=r^6-5r^4+19r^2+9,\qquad r\in\Q^\times.
\end{equation}
\end{corollary}
\begin{proof}
For $N=1$, the kernel has $v_1=0$, so the third offset vanishes. Each quartic row of Table~\ref{tab:normcatalog} permits only the already excluded fibers $r=0,\infty$. The four assignments with $N=r^2+9$ have the second square condition in~\eqref{eq:residual}. For $N=9r^2+1$, the second condition is $Z^2=9r^6+19r^4-5r^2+1$; substitution $r\mapsto1/r$ converts the pair to~\eqref{eq:residual} after rational-square rescalings. All statements about these eight norms are checked in the polynomial catalog.
\end{proof}

\section{The exact genus-five existence problem}\label{sec:genusfive}
We now prove Theorem~\ref{thm:main-d1}. The conic in~\eqref{eq:residual} has the parametrization
\begin{equation}\label{eq:conic}
 r=\frac{6h}{1-h^2},\qquad
 S=\frac{3(1+h^2)}{1-h^2},\qquad
 h=\frac{r}{S+3}.
\end{equation}
At finite nonzero $r$, the inverse has no exception, since $S=-3$ would force $r=0$.

\begin{lemma}\label{lem:pullback}
For $r=6h/(1-h^2)$,
\[
 (1-h^2)^6(r^6-5r^4+19r^2+9)=9P(h).
\]
Furthermore
\[
 P(h)=(h^6-13h^4+83h^2+1)(h^6+83h^4-13h^2+1)
\]
is squarefree. Its smooth double cover has genus five and contains exactly the eight points in~\eqref{eq:knownintro} above the parameter values $h=0,1,-1,\infty$.
\end{lemma}
\begin{proof}
The two polynomial identities follow by expansion. Squarefreeness follows from the explicit coprime factorization modulo $5$ in Appendix~\ref{app:d1checks}; exact Euclidean division independently gives $\gcd(P,P')=1$. A squarefree polynomial of degree twelve gives a double cover branched at twelve geometric points, so Riemann--Hurwitz yields $2g-2=-4+12$ and $g=5$. Substitution gives $P(0)=1$ and $P(1)=P(-1)=5184=72^2$. The infinity chart $t=1/h$, $\eta=Y/h^6$ has equation $\eta^2=P(t)$ and two smooth rational points $(0,\pm1)$, giving the two infinity points. This determines the fibers over the four indicated parameter values, not all rational points of the curve.
\end{proof}

\begin{proof}[Proof of Theorem~\ref{thm:main-d1}: forward direction]
An admissible $D1$ solution lies in one of the eight residual assignments by Corollary~\ref{cor:forty}. Apply the reciprocal substitution where needed, then~\eqref{eq:conic}. Lemma~\ref{lem:pullback} gives the rational point
\[
 (h,Y)=\left(\frac{r}{S+3},\frac{(1-h^2)^3Z}{3}\right)
\]
on $\mathcal C$. The original parameter $r$ is finite and nonzero; hence $h\notin\{0,1,-1,\infty\}$. The point lies outside $\mathcal K$.
\end{proof}

\begin{proof}[Proof of Theorem~\ref{thm:main-d1}: reconstruction]
Let $(h,Y)\in\mathcal C(\Q)\setminus\mathcal K$. Then $h$ is finite and $h\notin\{0,1,-1\}$. Put
\[
 Q(h)=h^4+34h^2+1,\qquad D(h)=(1+h^2)Q(h).
\]
Define three sine--cosine pairs by
\begin{align}
 (\sin\alpha,\cos\alpha)
 &=\left(\frac{1-h^2}{1+h^2},\frac{-2h}{1+h^2}\right),\notag\\
 (\sin\beta,\cos\beta)
 &=\left(\frac{12h(1-h^2)}{Q(h)},
           \frac{h^4-38h^2+1}{Q(h)}\right),\label{eq:reconstruction}\\
 (\sin\gamma,\cos\gamma)
 &=\left(\frac{-48h^2(1-h^2)}{D(h)},\frac{Y}{D(h)}\right).\notag
\end{align}
The first two pairs satisfy their unit-circle equations identically. The third does so because
\begin{equation}\label{eq:Didentity}
 D(h)^2-\bigl(48h^2(1-h^2)\bigr)^2=P(h).
\end{equation}
All denominators are positive for real $h$, so these are rational circle points.

With $L=(1-h^2)/D(h)$, their four signed offsets are
\begin{equation}\label{eq:reconstructedsines}
 L\bigl(h^4-62h^2+1,\ h^4-14h^2+1,\
         h^4+34h^2+1,\ -48h^2\bigr).
\end{equation}
They satisfy $s_0=s_1+s_3$ and $s_2=s_1-s_3$. It remains to prove nondegeneracy, since linear relations alone do not suffice.

Every entry, and every sum or difference of two entries in the parenthesized vector, factors into nonzero rational constants, powers of $h$, and polynomials of the form $h^4+ch^2+1$. Appendix~\ref{app:d1checks} lists all twelve pairwise expressions explicitly. The latter coefficients $c$ belong to
\[
 \{-110,-62,-38,-14,10,34,82\}.
\]
Each polynomial is monic with constant term one. Its only possible rational roots are $h=\pm1$, and its value there is $c+2\ne0$. Thus no entry vanishes and no two entries agree up to sign for the present rational $h$. The factor $L$ is also nonzero. Lemma~\ref{lem:roles} proves admissibility. The inverse integer change $x=\gamma-\beta$, $z=\alpha-\gamma+\beta$, $w=\beta$ gives an admissible point for $A_{D1}$. This proves the converse.
\end{proof}

\begin{corollary}\label{cor:fourtargets}
If $\mathcal C(\Q)=\mathcal K$, the configurations $D1,B2,C,E2$ have no admissible rational-circle point.
\end{corollary}
\begin{proof}
Theorem~\ref{thm:main-d1} excludes $D1$. The three determinant-two maps in~\eqref{eq:maps} and Lemma~\ref{lem:transport} exclude the other configurations.
\end{proof}

\begin{conjecture}\label{conj:points}
The eight-point list~\eqref{eq:knownintro} is all of $\mathcal C(\Q)$.
\end{conjecture}
By Theorem~\ref{thm:main-d1} this is exactly the statement that the configuration D1 — and, by Corollary~\ref{cor:fourtargets}, B2, C and E2 — satisfies the rational-circle rigidity conjecture of Section~\ref{sec:rigidity}. The completeness assertion is not used anywhere in this paper. For actual integer squares, the archive's three-prime note claims that no primitive square has a squarefree three-prime central root at all, which would cover the D1 case arithmetically; we have verified its structural lemmas and curves but not every class of its case analysis (Section~\ref{sec:prior}). Conjecture~\ref{conj:points} is the relaxed statement, independent of that claim. An additional rational point would instead give an explicit admissible $D1$ rational-circle point through~\eqref{eq:reconstruction}. It would not automatically give an integer-square magic square: the original Gaussian-prime conditions remain additional arithmetic restrictions.

\section{Quotients of the genus-five curve and the status of Conjecture~\ref{conj:points}}\label{sec:quotients}
This section was added on 22 October 2026. The polynomial $P$ in~\eqref{eq:Pintro} is even and palindromic, so $\mathcal C$ carries the involutions $\sigma_1:(h,Y)\mapsto(-h,Y)$ and $\sigma_2:(h,Y)\mapsto(1/h,Y/h^6)$.

\begin{proposition}[Quotients]\label{prop:quotients}
Put $g(s)=s^3+70s^2-1012s+6920$. Then
\[
 P(h)=f(h^2),\quad f(x)=x^6+70x^5-1009x^4+7060x^3-1009x^2+70x+1,\quad \frac{P(h)}{h^6}=g\!\left(h^2+h^{-2}\right),\quad \frac{f(x)}{x^3}=g\!\left(x+\frac1x\right).
\]
Hence $\mathcal C/\sigma_1$ is the genus-two curve $D:\ Y^2=f(x)$, which is bielliptic with quotients $E_\pm:\ v^2=(w\pm2)\,g(w)$, $w=x+1/x$; and $(h,Y)\mapsto(s,v)=(h^2+h^{-2},\,Y/h^3)$ maps $\mathcal C$ onto the elliptic curve $E_0:\ v^2=g(s)$. The three curves are isomorphic over $\Q$ to the LMFDB curves $123.\mathrm{a}1$ ($E_0$), $492.\mathrm{b}1$ ($E_+$) and $1968.\mathrm{j}1$ ($E_-$); the isomorphisms are verified exactly by the scaling of $c_4$ and $c_6$ with $u=4,8,12$ respectively.
\end{proposition}
\begin{proof}
The four identities are polynomial expansions, checked symbolically. The point $h=1$ of $\mathcal K$ maps to $(s,v)=(2,72)$ on $E_0$. The identifications are by $j$-invariant followed by the exact $c_4,c_6$ test, which distinguishes a curve from its twists.
\end{proof}

\begin{proposition}[Positive rank]\label{prop:posrank}
Each of $E_0,E_+,E_-$ has a rational point of infinite order. Hence $\operatorname{rank}E_i(\Q)\ge1$ for all three, and $\operatorname{rank}\operatorname{Jac}(\mathcal C)(\Q)\ge3$.
\end{proposition}
\begin{proof}
On the integral model $v^2=g(s)$ the point $(2,72)$ has $3\cdot(2,72)$ with abscissa $11042/289$; by the Nagell--Lutz theorem a torsion point of an integral Weierstrass model has integral coordinates, so $(2,72)$ has infinite order. On the integral monic models of $E_+$ and $E_-$ obtained by moving the rational root $w=\mp2$ of the quartic to infinity, the points $(0,9216)$ and $(0,5184)$ have the non-integral multiples $8P$ (abscissa $-2304/25$) and $3P$ (abscissa $113291136/292681$). The three curves have pairwise distinct $j$-invariants and conductors $123$, $492$, $1968$, so they are pairwise non-isogenous over $\Q$; their pullbacks are independent abelian subvarieties of $\operatorname{Jac}(\mathcal C)$, and the ranks add.
\end{proof}

The LMFDB records rank exactly one for each curve (a database fact we quote and do not re-prove), so $\operatorname{rank}\operatorname{Jac}(D)(\Q)=2=g(D)$. Consequently neither the rank-zero descents of Section~\ref{sec:norm} nor classical Chabauty on $D$ can prove Conjecture~\ref{conj:points}; the conjecture is an instance of quadratic Chabauty for a bielliptic genus-two curve of Mordell--Weil rank two, with the generators $(-7,18)$ on $492.\mathrm{b}1$ and $(-2,3)$ on $1968.\mathrm{j}1$ as inputs, or of elliptic-curve Chabauty over a quadratic field. We record also that a sieved exact search finds no rational point of $\mathcal C$ with $\max(|p|,q)\le10^6$ outside $\mathcal K$, and none on $D$ with height at most $3\cdot10^5$ outside the images of $\mathcal K$; this is evidence, not proof.

\section{A conjecture: rational-circle rigidity}\label{sec:rigidity}

Everything in this paper concerns the \emph{rational-circle relaxation}: an integer frequency matrix $A\in\Z^{4\times r}$
satisfying the \emph{full-depth condition} that in every column at least three rows attain the column's maximal
absolute value (the condition Lemma~\ref{lem:depth} forces for every primitive square; for a squarefree center it says
that every coordinate is nonzero in at least three rows of a matrix with entries in $\{0,\pm1\}$), together with the
question of whether a point of $T(\Q)^r$ is admissible for $A$ in the sense of Section~\ref{sec:introduction}. Two equations in $r$ parameters: the solution set has
dimension $r-2$, finite for two parameters, curves for three, and surfaces for $r\ge4$. The results above suggest a
single statement.

\begin{conjecture}[Rational-circle rigidity]\label{conj:rigidity}
For every $r\ge2$ and every full-depth configuration $A$, there is no admissible point of $T(\Q)^r$ for $A$.
\end{conjecture}

The conjecture is strictly stronger than the nonexistence of a primitive magic square of nine distinct squares: every
such square yields a full-depth configuration and an admissible point of $T(\Q)^r$ (Lemma~\ref{lem:depth} and~\eqref{eq:areaangle}),
and the conjecture discards the Gaussian-prime condition entirely. Theorem~\ref{thm:main-two} proves it for two parameters
through exponent sum seven and Corollary~\ref{cor:parallelogram} for every parallelogram configuration, over the reals; for two parameters at all exponents the conjecture would give a purely trigonometric proof of the archive's theorem
that no primitive square has a two-prime central root (Section~\ref{sec:prior}), a theorem presently proved only through the
Gaussian-prime arithmetic;
Conjecture~\ref{conj:points} is, by Theorem~\ref{thm:main-d1}, equivalent to its instance for the configuration D1, and by
Corollary~\ref{cor:fourtargets} implies its instances for B2, C and E2. If the conjecture fails, the smallest $r$ at which it
fails locates where the arithmetic of the circle points becomes essential: below that $r$ the obstruction to magic squares
is trigonometric, and from that $r$ on it is arithmetic. The three-parameter cases being open, the first unexplored test is
$r=4$, where the solution sets are surfaces. The two results of this section concern that case.

\subsection{Four parameters at bounded height}

Write $\tan(\chi_j/2)=[p_j:q_j]\in\PP(\Q)$ in lowest terms, with height $\max(|p_j|,|q_j|)$; thus $0=[0:1]$ and
$\infty=[1:0]$ have height one, and correspond to $\chi_j=0$ and $\chi_j=\pi$.

\begin{proposition}\label{prop:four}
The $11{,}862$ four-element sets of nonzero vectors of $\{0,\pm1\}^4$ modulo sign that satisfy the full-depth condition
form $113$ orbits under signed coordinate permutations. No such configuration admits an admissible point of $T(\Q)^4$ at
which three of the four parameters $\tan(\chi_j/2)$ have projective height at most $4$.
\end{proposition}

\begin{proof}[Computer-assisted proof]
For each orbit representative, each choice of a coordinate $j$, and each triple of values of the other parameters among the
$24$ points of $\PP(\Q)$ of height at most $4$ (including $0$ and $\infty$), every row's sine is an integer linear form in
$(\cos\chi_j,\sin\chi_j)$ once the positive common denominator is cleared. Each of the $48$ signed systems of
Lemma~\ref{lem:roles} is two such linear equations, and three cases arise. If their determinant is nonzero, the unique
solution is tested exactly against the unit circle. If the determinant vanishes but the equations are not both identically
zero, the consistent cases reduce to a line meeting the circle, decided by an integer square-root test. If both equations
vanish identically — which happens at $12{,}240$ of the $113\cdot4\cdot13{,}824\cdot48$ system evaluations, at parameter
triples containing torsion values — the relations hold for \emph{every} value of $\chi_j$, and the four sines
$A_rc+B_rs+C_r$ form a one-parameter family; such a family contains admissible points (all but finitely many) if and only if
no coefficient triple $(A_r,B_r,C_r)$ vanishes and no two are equal up to sign, which is tested exactly, and in every
occurrence the family is degenerate. Surviving points are tested for admissibility exactly. Since every admissible point has
some three parameters in the grid for some ordering of the coordinates, and every configuration is a signed permutation of a
representative, the sweep is complete for the stated height. It finds no admissible point; height $3$ was swept separately
with the same outcome, and an earlier sweep of the finite chart at height $5$ (without $0,\infty$ and without the family
branch) also found none. The code (\path{verification/four_prime/r4search.py}), its checkpoints, and unit tests —
reconstruction of all four sines from the linear forms at $400$ random rational points, recovery of planted circle points,
and the admissibility and singular branches — accompany the paper.
\end{proof}

A heuristic count predicts the opposite: a surface in four rational parameters should carry many rational points. The
next result explains the absence for the most symmetric configuration.

\subsection{The Hadamard solution surfaces are of general type}

Let $A_{\mathrm H}$ be the Hadamard configuration with rows $(1,1,1,1),(1,1,-1,-1),(1,-1,1,-1),(1,-1,-1,1)$, and put
\[
\alpha=\chi_1+\chi_2,\quad \beta=\chi_3+\chi_4,\quad \gamma=\chi_1-\chi_2,\quad \delta=\chi_3-\chi_4,\qquad
u=\gamma+\delta,\quad v=\gamma-\delta,
\]
so that $(s_0,s_1,s_2,s_3)=(\sin(\alpha+\beta),\sin(\alpha-\beta),\sin(\gamma+\delta),\sin(\gamma-\delta))$. In every signed
system~\eqref{eq:roles} the two equations can be solved for $(s_2,s_3)$ as linear forms in $(s_0,s_1)$ — the coefficient
block on $(s_2,s_3)$ is invertible in all three index patterns — so each system reads
\begin{equation}\label{eq:bidouble}
\sin u=F_1(\alpha,\beta),\qquad \sin v=F_2(\alpha,\beta),
\end{equation}
with $F_1,F_2\in\Q\sin(\alpha+\beta)+\Q\sin(\alpha-\beta)$. Up to the symmetries $u\mapsto-u$, $v\mapsto-v$,
$u\leftrightarrow v$ and $\alpha\leftrightarrow\beta$, exactly four pairs $(F_1,F_2)$ occur:
\[
\begin{array}{cll}
\text{(I)} & \sin\alpha\cos\beta, & \cos\alpha\sin\beta;\\
\text{(II)} & 2\cos\alpha\sin\beta, & 2\sin\alpha\cos\beta;\\
\text{(III)} & 2\cos\alpha\sin\beta, & 3\cos\alpha\sin\beta-\sin\alpha\cos\beta;\\
\text{(IV)} & 2\cos\alpha\sin\beta, & \sin\alpha\cos\beta+3\cos\alpha\sin\beta .
\end{array}
\]

\begin{theorem}[All Hadamard solution surfaces have general type]\label{thm:gentype}
For each of the $48$ signed systems of $A_{\mathrm H}$, let $S\subset(\PP)^4$ be the closure of its solution set in the
coordinates $\tan(\chi_j/2)$, and let $S''\subset(\PP)^4$ be the closure of~\eqref{eq:bidouble} in the coordinates
$\tan(\alpha/2),\tan(\beta/2),\tan(u/2),\tan(v/2)$. Then:
\begin{enumerate}
\item $S''$ is the irreducible $(\Z/2)^2$-cover of $\PP\times\PP$ branched along $B_1=\{F_1=\pm1\}$ and
$B_2=\{F_2=\pm1\}$, each of bidegree $(4,4)$, reduced, with no common component. Every component of $B_1\cup B_2$ has only
ordinary nodes as singularities, the two components of each $B_i$ meet transversally, no three components pass through
one point, the nodes of $B_1$ lie off $B_2$ and conversely, and every point of $B_1\cap B_2$ is either transversal or a
simple tangency.
\item Consequently $S''$ has only rational double points of type $A_1$, and its minimal resolution is a minimal surface of
general type with
\[
p_g=11,\qquad q=0,\qquad K^2=32,\qquad \chi(\mathcal O)=12 .
\]
\item The change of angles is an isogeny of tori of degree $8$ and induces a dominant, generically finite rational map
$S\dashrightarrow S''$ of degree $8$ (the induced map of the toric compactifications need not be a morphism), under which
every two-dimensional component of $S$ dominates $S''$; hence every such component has Kodaira dimension~$2$.
\end{enumerate}
\end{theorem}

\begin{proof}
(1) In half-angle coordinates $(a,b,w,y)$, \eqref{eq:bidouble} becomes two equations of multidegrees $(2,2,2,0)$ and
$(2,2,0,2)$; the first is the double cover $F_1w^2-2w+F_1=0$ of the $(a,b)$-plane, branched where its discriminant
$4(1-F_1^2)$ vanishes, i.e.\ on $B_1$, and likewise for $B_2$. Each $\{F_i=\pm1\}$ clears to a polynomial of bidegree
$(2,2)$. The fibre product of two irreducible double covers with distinct branch loci is irreducible. The geometry of the
branch curves is a finite Gr\"obner-basis computation in all four affine charts of $\PP\times\PP$
(\path{verification/four_prime/branchcheck.py}, \path{branchfine.py}), with the following outcome. Type~(I): each $(2,2)$
component splits over $\Q(i)$ into two smooth $(1,1)$ curves, e.g.\ $(1+a^2)(1+b^2)-2a(1-b^2)=(1-a)^2+b^2(1+a)^2$, so
$B_i$ is four rational curves with $2+2+8=12$ transversal nodes, and $B_1\cap B_2$ consists of $32$ transversal points.
Type~(II): all four components are smooth curves of genus one; $B_1$ and $B_2$ each have $8$ transversal nodes, and
$B_1\cap B_2$ consists of $16$ simple tangencies and no transversal points — the full intersection number
$(4,4)\cdot(4,4)=32$ — at $\tan(\alpha/2),\tan(\beta/2)\in\{\pm1\pm\sqrt2\}$, i.e.\ at $\alpha,\beta\in\{\pm\pi/4,\pm3\pi/4\}$,
where $\cos2\alpha=\cos2\beta=0$. Types~(III), (IV): the components of $B_1$
are smooth of genus one, the components of $B_2$ split over $\Q(\sqrt2)$ into smooth $(1,1)$ curves (the coefficients
$3\pm2\sqrt2=(1\pm\sqrt2)^2$ appear), all intersections within each $B_i$ are transversal nodes, the nodes lie off the other
divisor, there are no triple points, and $B_1\cap B_2$ consists of $16$ transversal points and $8$ simple tangencies
($16+2\cdot8=32$). In each type the counts were confirmed against the intersection number.

(2) Local analysis of the cover $z^2=f$, $w^2=g$ over a point of the base: over a smooth point of one branch curve not on
the other it is smooth; over a node of $B_1$ (resp.\ $B_2$) it is $z^2=xy$ on each of two sheets, an $A_1$ point; over a
transversal point of $B_1\cap B_2$ it is $z^2=x$, $w^2=y$, smooth; and over a simple tangency, with $f=y$, $g=y-x^2$, it is
$w^2-z^2+x^2=0$, again an $A_1$ point. Hence all singularities of $S''$ are canonical. The cover is a normal bidouble cover
with branch data $D_1=B_1$, $D_2=B_2$, $D_3=0$; by Pardini's structure theorem for normal abelian covers~\cite{Pardini},
$\pi_*\mathcal O_{S''}=\bigoplus_i\mathcal O(-L_i)$ with $2L_1=D_2$, $2L_2=D_1$, $2L_3=D_1+D_2$, so $L_1=L_2=(2,2)$ and
$L_3=(4,4)$. Rational double points are rational and crepant, so the minimal resolution has
\[
p_g=\sum_{i=1}^{3}h^0\bigl(K_X+L_i\bigr)=h^0(\mathcal O)+h^0(\mathcal O)+h^0(\mathcal O(2,2))=1+1+9=11,\qquad
q=\sum_{i=1}^{3}h^1(K_X+L_i)=0,
\]
\[
K^2=4\bigl(K_X+\tfrac12(D_1+D_2)\bigr)^2=4\,(2,2)^2=32 .
\]
Independently, adjunction on the complete intersection gives $K_{S''}=\mathcal O(2,2,0,0)|_{S''}$, the pullback of an ample class
under a finite map and hence ample, with $K^2=(2h_1+2h_2)^2(2h_1+2h_2+2h_3)(2h_1+2h_2+2h_4)=32$, and the Koszul resolution
gives $\chi(\mathcal O_{S''})=1-(-1)-(-1)+9=12=1-q+p_g$. Ampleness of $K$ on the singular model shows the resolution is
minimal; Noether's inequality $2p_g-4\le K^2\le 9\chi$ reads $18\le32\le108$~\cite{BHPV}.

(3) The change of angles $\chi\mapsto(\alpha,\beta,\gamma,\delta)$ has determinant $4$ and $(\gamma,\delta)\mapsto(u,v)$
has determinant $2$; on the tori this is an isogeny of degree $8$, so the induced rational map $S\dashrightarrow S''$ is
dominant and generically finite of degree $8$, and every surface component of $S$ dominates the irreducible surface $S''$.
Kodaira dimension does not decrease under a dominant generically finite rational map (resolve the indeterminacy and apply
the statement for morphisms).
\end{proof}

\begin{corollary}
Under the Bombieri--Lang conjecture~\cite{Lang}, every component of every Hadamard solution surface has only finitely many
rational points outside finitely many curves. This is consistent with Proposition~\ref{prop:four} and with
Conjecture~\ref{conj:rigidity}; it proves neither.
\end{corollary}

\subsection{Two further hypotheses}

Conjecture~\ref{conj:rigidity} asserts emptiness, which is far stronger than what the geometry suggests: for a surface or
higher-dimensional variety of general type, the Bombieri--Lang conjecture predicts only that rational points are not
Zariski dense. We therefore record the statement that the geometry supports directly, and a structural hypothesis about
where the rational points that do exist are located.

\begin{conjecture}[Weak form]\label{conj:weak}
For every full-depth configuration, the admissible points of $T(\Q)^r$ are not Zariski dense in any component of the
solution variety. By Bombieri--Lang this holds for every component of general type, in particular for every component in
Theorem~\ref{thm:gentype}.
\end{conjecture}

Every rational point of a solution variety found in this paper lies on a \emph{torsion coset}: the eight points of $K$
on $\mathcal C$ have $\alpha,\beta,\gamma\in\{0,\pm\pi/2,\pi\}$; the trivial sieve classes in the extended record are
torsion points; and all $12{,}240$ free-angle families in the height-$4$ sweep occur at parameter triples containing a
torsion value. The torsion of $T(\Q)$ is $\{\pm1,\pm i\}$, and a rational point of $T(\Q)^r$ whose coordinates are
multiplicatively independent lies on no proper subtorus coset, so the following is a genuine restriction.

\begin{conjecture}[Torsion-coset hypothesis]\label{conj:torsion}
For every full-depth configuration, every rational-circle point of its solution variety satisfies a nontrivial relation
$\prod_j z_j^{\,n_j}\in\{\pm1,\pm i\}$ with $(n_j)\in\Z^r\setminus\{0\}$; that is, the rational points lie on
finitely many torsion cosets of proper subtori.
\end{conjecture}

The hypothesis has an inductive consequence. If the relation has some coefficient $n_j=\pm1$, eliminating $\chi_j$
turns the configuration into one in $r-1$ parameters, with integer entries of larger absolute value in general, so
Conjecture~\ref{conj:rigidity} for $r$ would follow from Conjecture~\ref{conj:torsion} for $r$ together with
Conjecture~\ref{conj:rigidity} for $r-1$; relations with all $|n_j|\ge2$ instead pass to a finite cover. In this sense
the torsion-coset hypothesis would reduce rigidity in all dimensions to the three-parameter cases, of which D1 — the
completeness of $K$ — is the first. It is also the mechanism by which a search finds nothing: the rational points that
the equations do admit sit on degenerate loci where two sine magnitudes coincide identically. We offer it as the
hypothesis whose confirmation or refutation would most change the picture.

\subsection{Evidence for the torsion-coset hypothesis and a degeneracy-rank lemma (22 October 2026)}\label{sec:torsionprobe}
The sweep of Proposition~\ref{prop:four} was re-instrumented to record \emph{every} exact solution of the $48$ signed systems, admissible or not, for all $113$ orbits, all four choices of solved coordinate, and all triples of the other parameters of projective height at most $H$, and each solution was tested for Conjecture~\ref{conj:torsion} by a complete criterion: with $\tan(\chi_j/2)=p_j/q_j$ one has $z_j=(q_j+ip_j)/(q_j-ip_j)$, whose image in $\Q(i)^\times/\{\pm1,\pm i\}$ is determined by the Gaussian-prime factorisation of $q_j+ip_j$; a relation $\prod z_j^{n_j}\in\{\pm1,\pm i\}$ exists if and only if the four exponent vectors are linearly dependent, and the rank of these vectors (the \emph{exponent rank}) is computed exactly.

\begin{proposition}\label{prop:torsiondata}
At $H=3$ there are $273{,}024$ distinct rational points of the solution varieties, at $H=4$ there are $662{,}528$, and at $H=5$ there are $1{,}916{,}672$ (each grid contains the previous one). Every one of them satisfies a nontrivial relation $\prod z_j^{n_j}\in\{\pm1,\pm i\}$; none is admissible; every one has a vanishing sine and two sines of equal magnitude. The exponent rank is at most $2$ at $H=3$; at $H=4$ it equals $3$ for $13{,}824$ points and at $H=5$ for $235{,}008$ points, in the same two orbits.
\end{proposition}
The rank-three points refute the sharper statement, suggested by the $H\le3$ data, that every rational point lies on a torsion coset of a subtorus of dimension at most two. They are explained by the following elementary lemma.

\begin{lemma}[Degeneracy relations]\label{lem:degrel}
Let $z\in T(\Q)^r$ and $s_j=\sin(A_j\cdot\chi)$. If $s_r=0$ then $\prod_j z_j^{A_{rj}}=\pm1$. If $|s_i|=|s_k|$ with $i\ne k$ then $\prod_j z_j^{A_{ij}-A_{kj}}=\pm1$ or $\prod_j z_j^{A_{ij}+A_{kj}}=\pm1$, with a nonzero exponent vector because the rows are distinct projective vectors. Consequently every rational point with a vanishing sine or two sines of equal magnitude satisfies Conjecture~\ref{conj:torsion}; a point with both degeneracies has exponent rank at most $3$, with equality only if $A_i\pm A_k$ is a rational multiple of the vanishing row $A_r$.
\end{lemma}
\begin{proof}
$s_r=0$ means $A_r\cdot\chi\in\pi\Z$, and $|s_i|=|s_k|$ means $A_i\cdot\chi\equiv\pm A_k\cdot\chi\pmod\pi$; exponentiate. Two independent relations lower the rank to at most $2$.
\end{proof}
\begin{theorem}[Trichotomy on a signed system]\label{thm:trichotomy}
Let $z\in T(\Q)^r$ satisfy a signed system~\eqref{eq:roles}, $s_i=\varepsilon s_k+\eta s_l$, $s_j=\varepsilon s_k-\eta s_l$, and put $A=\max(|s_k|,|s_l|)$, $B=\min(|s_k|,|s_l|)$. Then $\{|s_i|,|s_j|\}=\{A+B,A-B\}$, and exactly one of the following holds: (i) some sine vanishes; (ii) $z$ is admissible; (iii) $A=2B>0$. In case (i) two of the remaining sines have equal magnitude, so by Lemma~\ref{lem:degrel} the point lies on a torsion coset and has exponent rank at most $3$.
\end{theorem}
\begin{proof}
For real $x,y$ the magnitudes of $x+y$ and $x-y$ are $A+B$ and $A-B$ with $A=\max(|x|,|y|)$, $B=\min(|x|,|y|)$. If no sine vanishes then $B>0$ and $A\ne B$, and the four magnitudes $A,B,A+B,A-B$ are pairwise distinct unless $A-B=B$; by Lemma~\ref{lem:roles} the point is admissible exactly when they are distinct. For (i): if $s_l=0$ then $s_i=s_j$; if $s_k=0$ then $|s_i|=|s_j|=|s_l|$; if $s_i=0$ or $s_j=0$ then $|s_k|=|s_l|$.
\end{proof}
Case (iii) is the only way a rational point of a solution variety can avoid both admissibility and a monomial relation, and it is not a torsion-coset condition. Every one of the $1{,}916{,}672$ points recorded at $H\le5$ (each grid contains the previous one, so the sweeps are not added) is of type (i); a direct classification of every point found no point of type (iii). We therefore state the hypothesis that the data are consistent with and that implies both Conjecture~\ref{conj:rigidity} and Conjecture~\ref{conj:torsion}; it is strictly stronger than Conjecture~\ref{conj:rigidity}, since it also excludes type (iii), and the evidence for the extra exclusion is only the absence of type (iii) points at $H\le5$, where the grid carries the split primes $5$, $13$, $17$, $29$ and $41$:
\begin{conjecture}[Vanishing-sine rigidity]\label{conj:vanishing}
For every full-depth configuration, every rational point of every signed system has a vanishing sine.
\end{conjecture}
Conjecture~\ref{conj:vanishing} is equivalent to Conjecture~\ref{conj:rigidity} together with the nonexistence of rational points of type (iii). The next computation is therefore not a larger admissibility sweep but a direct search for type (iii) points: rational-circle solutions of a signed system with $|s_l|=2|s_k|$ or $|s_k|=2|s_l|$ and all sines nonzero.

The parallel-triple condition of the lemma holds for exactly $23$ of the $113$ orbit representatives, and the two orbits that carry rank-three points at $H=4$ and at $H=5$ are among them; every recorded rank-three point has a vanishing row parallel to the sum or difference of its two equal-magnitude rows. The remaining $21$ orbits are candidates for larger heights; the criterion is necessary, not sufficient. The lemma also shows that Conjecture~\ref{conj:rigidity} implies Conjecture~\ref{conj:torsion} for every solution-variety point that fails admissibility through a vanishing sine or coinciding magnitudes; the only other failure mode, $a=2b$, yields no monomial relation and was never observed. The code, the per-orbit counts and the twenty recorded rank-three examples are in the companion repository.

Theorem~\ref{thm:gentype} shows that the general-type behaviour Bremner~\cite{Bremner2001} and
Bruin--Thomas--V\'arilly-Alvarado~\cite{BTV} established for the magic-square surface is already present in the trigonometric
relaxation, with the Gaussian-prime arithmetic removed. The decoupling that makes the proof elementary — the rows pair into
$\{v_0,v_1\}$ and $\{v_2,v_3\}$ with $v_0\pm v_1$ and $v_2\pm v_3$ supported on complementary coordinate pairs — is special
to the Hadamard configuration. For a general full-depth configuration the two relations are divisors of multidegree $(2,2,2,2)$ in $(\PP)^4$; whenever
they have no common component the solution surface is their complete intersection, adjunction gives
$K_S=\mathcal O(2,2,2,2)|_S$, which is ample, and the surface is of general type whenever its singularities are canonical. Verifying that for the other $112$ orbits, and
locating any configuration whose surface is instead rational — which would be a direct route to a rational-circle solution
and would refute Conjecture~\ref{conj:rigidity} — is the natural continuation.


\section{Proof data and reproducibility}\label{sec:repro}
\emph{Companion repository (22 October 2026).} The additions of this version were produced and checked by a multi-agent pipeline that exposes the verification tree of this paper as tools (an MCP server), attacks the open problems by fixed playbooks, submits every claim to an adversarial referee agent, and absorbs accepted results as additional stages of \texttt{verify\_paper.py} without modifying the shipped proof records. The repository \url{https://github.com/ksamiein-eng/magic-squares-agents} contains the certificates for exponent sums eight and nine, the quotient-curve computations, the torsion-coset probe, the degeneracy-rank script, verbatim agent transcripts, a claim map, provenance hashes, and the extended suite \texttt{verify\_paper\_plus.py} (the sixteen stages of this paper followed by six agent-added stages, all passing).

The principal results have different computational dependencies. Theorem~\ref{thm:main-two} uses a finite elimination certificate; Theorem~\ref{thm:main-d1} uses a forty-eight-row polynomial catalog, three finite descent calculations, and explicit identities. Their logical dependencies are summarized in Table~\ref{tab:proofobjects}.
\begin{table}[htbp]
\centering\small
\begin{tabular}{>{\raggedright\arraybackslash}p{0.32\linewidth}>{\raggedright\arraybackslash}p{0.60\linewidth}}
\toprule
Result & Certificate or independent check\\
\midrule
Theorem~\ref{thm:main-two} & \path{verification/two_prime/results/certificate_degree7.json}; repeated-addition and Chebyshev/Sylvester checkers.\\[3pt]
Lemma~\ref{lem:normcatalog} & \path{verification/d1_norm/d1_certificate.json}; polynomial kernels, divisors, both norms, and all assignment tuples.\\[3pt]
Proposition~\ref{prop:quartics} & The same D1 record; exhaustive descent-class obstructions and good-prime counts.\\[3pt]
Theorem~\ref{thm:main-d1} & Exact conic pullback, squarefreeness, circle reconstruction, and nondegeneracy identities in the D1 checker and \path{check_paper.py}.\\[3pt]
Proposition~\ref{prop:orbits}; \eqref{eq:maps} & Independent finite enumeration and exact signed matrix products in \path{check_paper.py}.\\[3pt]
Proposition~\ref{prop:four} & \path{verification/four_prime/r4search.py}; complete projective sweeps at heights $3$ and $4$ (\path{results/r4_H3_final.json}, \path{r4_H4_final.json}) with exact handling of the Cramer, line--circle and free-angle branches; unit tests of the linear forms, Cramer solver, admissibility and singular branch. The legacy finite-chart height-5 sweep is \path{r4_H5_affine_legacy.json}.\\[3pt]
Theorem~\ref{thm:gentype} & \path{verification/four_prime/branchcheck.py}, \path{branchfine.py}, \path{general_type_check.py}: branch bidegrees, splitting fields, node and tangency classification in all charts, triple-point and common-component exclusion, and the two invariant computations.\\
\bottomrule
\end{tabular}
\caption{Proof objects for the claims selected for this article.}\label{tab:proofobjects}
\end{table}

The companion archive contains the source manuscript, bibliography, certificates, checker source, and execution records. The main command is
\begin{verbatim}
python -m pip install -r requirements.txt
python verify_paper.py
python verification/four_prime/general_type_check.py
python verification/archive_audit/archive_checks.py   # curves and identities of the archive audit
python verification/four_prime/branchfine.py 2   # likewise 4 and 6
\end{verbatim}
The four-parameter sweep is reproduced with \path{python verification/four_prime/r4search.py 4} (about fifteen minutes on one core; resumable by orbit range \path{START END}).
The Python standard library suffices for the two-prime acceptance path and the D1 arithmetic checker. The separate symbolic checks use SymPy 1.14.0~\cite{SymPy}. Verification runs in temporary copies, and assertions must remain enabled. The producer uses factorization to discover the two-prime certificate; acceptance checks the displayed identities and modular witnesses rather than trusting the producer's conclusions.

The use of two implementations reduces shared implementation risks but does not replace review of the mathematical reductions. Python, the checker implementations, and SymPy's exact polynomial routines used by the symbolic checks are within the software trust boundary. No formal proof-assistant verification is claimed. The principal theorems require neither published positive-rank generator data nor an executed Sage, PARI, Magma, or Chabauty calculation. The classical descent and torsion-reduction results used in Proposition~\ref{prop:quartics} are cited explicitly.

\paragraph{Use of generative AI.}
ChatGPT and Claude were used for symbolic and computational work, implementation of verification code, and preparation and revision of the exposition. The additions dated 22 October 2026 (Theorem~\ref{thm:sumnine}, Section~\ref{sec:quotients}, Section~\ref{sec:torsionprobe}) were generated, machine-verified and refereed by Claude agents orchestrated through Claude Code; the authors reviewed the statements and take responsibility for all claims.

\appendix
\section{Two-parameter coverage and elimination certificates}\label{app:representatives}
This appendix separates three obligations in Theorem~\ref{thm:main-two}: covering the frequency configurations, checking the exact polynomial eliminations, and excluding their remaining rational fibers. A count of successful systems addresses only the second obligation unless the first and third are supplied as well.

\subsection{Reduction to coordinate-primitive configurations}
For a normal form $(k,\ell;m,n)$, set
\[
 d=\gcd(k,|m|),\qquad e=\gcd(\ell,|n|),\qquad
 (K,L;M,N)=\left(\frac{k}{d},\frac{\ell}{e};\frac{m}{d},\frac{n}{e}\right),
\]
with $\gcd(k,0)=k$. Then
\[
 A_{k,\ell;m,n}=A_{K,L;M,N}\diag(d,e).
\]
An admissible point for the left-hand side gives one for the reduced matrix at $(z_1^d,z_2^e)$. This implication is sufficient; no rational inverse of the power map is used.

After reduction, interchange columns if $K<L$; this also interchanges $M,N$, up to row signs and ordering. Simultaneous replacement $(M,N)\mapsto(-M,-N)$ gives the remaining sign identification. When $K=L$, column interchange is an additional symmetry. If $M=N=0$, the four frequencies form a parallelogram and Corollary~\ref{cor:parallelogram} applies. Enumerating the other reduced configurations through $K+L\le7$ gives the 43 representatives below. Independently reducing every raw form by these rules accounts for all 168 forms, including those with a zero deficient coordinate.

\begin{center}
\small\setlength{\tabcolsep}{8pt}\renewcommand{\arraystretch}{1.15}
\begin{tabular}{crrr}
\toprule
$(k,\ell)$ & Raw forms & Direct cases & Signed systems\\
\midrule
$(1,1)$ & 1 & 0 & 0\\
$(2,1)$ & 3 & 1 & 48\\
$(3,1)$ & 5 & 2 & 96\\
$(2,2)$ & 9 & 2 & 96\\
$(4,1)$ & 7 & 2 & 96\\
$(3,2)$ & 15 & 4 & 192\\
$(5,1)$ & 9 & 4 & 192\\
$(4,2)$ & 21 & 4 & 192\\
$(3,3)$ & 25 & 6 & 288\\
$(6,1)$ & 11 & 2 & 96\\
$(5,2)$ & 27 & 8 & 384\\
$(4,3)$ & 35 & 8 & 384\\
\midrule
Total & 168 & 43 & 2,064\\
\bottomrule
\end{tabular}
\end{center}


Here ``raw forms'' means the $(2k-1)(2\ell-1)$ choices of $(m,n)$ for the given bounds, before any identifications. A direct case has 48 signed systems. Thus 168 is a configuration count, 43 a representative count, and 2,064 an equation-pair count; these are not interchangeable.

{\small\setlength{\tabcolsep}{6pt}\renewcommand{\arraystretch}{1.12}
\begin{longtable}{ccccr}
\caption{The 43 direct cases. Every row contains 48 signed systems. $d_R$ is the largest actual resultant degree, not the larger degree bound used for verification. The last column counts distinct nonlinear no-root factors within the row.}\label{tab:representatives}\\
\toprule
$(k,\ell)$ & $(m,n)$ & $d_R$ & Possible rational $t$ & No-root factors\\
\midrule\endfirsthead
\toprule
$(k,\ell)$ & $(m,n)$ & $d_R$ & Possible rational $t$ & No-root factors\\
\midrule\endhead
$(2,1)$ & $(1,0)$ & 14 & $0,\pm1$ & 25\\
$(3,1)$ & $(1,0)$ & 22 & $0,\pm1$ & 26\\
$(3,1)$ & $(2,0)$ & 22 & $0,\pm1$ & 29\\
$(2,2)$ & $(1,1)$ & 30 & $0$ & 35\\
$(2,2)$ & $(1,-1)$ & 30 & $0$ & 25\\
$(4,1)$ & $(1,0)$ & 30 & $0,\pm1$ & 27\\
$(4,1)$ & $(3,0)$ & 30 & $0,\pm1$ & 27\\
$(3,2)$ & $(1,1)$ & 46 & $0,\pm1$ & 29\\
$(3,2)$ & $(1,-1)$ & 46 & $0,\pm1$ & 32\\
$(3,2)$ & $(2,1)$ & 46 & $0,\pm1$ & 31\\
$(3,2)$ & $(2,-1)$ & 46 & $0,\pm1$ & 31\\
$(5,1)$ & $(1,0)$ & 38 & $0,\pm1$ & 27\\
$(5,1)$ & $(2,0)$ & 38 & $0,\pm1$ & 31\\
$(5,1)$ & $(3,0)$ & 38 & $0,\pm1$ & 29\\
$(5,1)$ & $(4,0)$ & 38 & $0,\pm1$ & 31\\
$(4,2)$ & $(1,1)$ & 62 & $0$ & 27\\
$(4,2)$ & $(1,-1)$ & 62 & $0$ & 27\\
$(4,2)$ & $(3,1)$ & 62 & $0$ & 27\\
$(4,2)$ & $(3,-1)$ & 62 & $0$ & 31\\
$(3,3)$ & $(1,1)$ & 70 & $0,\pm1$ & 60\\
$(3,3)$ & $(1,-1)$ & 70 & $0,\pm1$ & 53\\
$(3,3)$ & $(1,2)$ & 70 & $0,\pm1$ & 25\\
$(3,3)$ & $(1,-2)$ & 70 & $0,\pm1$ & 28\\
$(3,3)$ & $(2,2)$ & 70 & $0,\pm1$ & 37\\
$(3,3)$ & $(2,-2)$ & 70 & $0,\pm1$ & 27\\
$(6,1)$ & $(1,0)$ & 46 & $0,\pm1$ & 27\\
$(6,1)$ & $(5,0)$ & 46 & $0,\pm1$ & 29\\
$(5,2)$ & $(1,1)$ & 78 & $0,\pm1$ & 29\\
$(5,2)$ & $(1,-1)$ & 78 & $0,\pm1$ & 33\\
$(5,2)$ & $(2,1)$ & 78 & $0,\pm1$ & 27\\
$(5,2)$ & $(2,-1)$ & 78 & $0,\pm1$ & 29\\
$(5,2)$ & $(3,1)$ & 78 & $0,\pm1$ & 31\\
$(5,2)$ & $(3,-1)$ & 78 & $0,\pm1$ & 32\\
$(5,2)$ & $(4,1)$ & 78 & $0,\pm1$ & 31\\
$(5,2)$ & $(4,-1)$ & 78 & $0,\pm1$ & 31\\
$(4,3)$ & $(1,1)$ & 94 & $0,\pm1$ & 53\\
$(4,3)$ & $(1,-1)$ & 94 & $0,\pm1$ & 55\\
$(4,3)$ & $(1,2)$ & 94 & $0$ & 29\\
$(4,3)$ & $(1,-2)$ & 94 & $0$ & 29\\
$(4,3)$ & $(3,1)$ & 94 & $0,\pm1$ & 57\\
$(4,3)$ & $(3,-1)$ & 94 & $0,\pm1$ & 55\\
$(4,3)$ & $(3,2)$ & 94 & $0$ & 29\\
$(4,3)$ & $(3,-2)$ & 94 & $0$ & 29\\
\bottomrule
\end{longtable}
}


\subsection{Reading a certificate record}
The file \path{verification/two_prime/results/certificate_degree7.json} has one record for each direct case. Its \texttt{bounds}, \texttt{drops}, and \texttt{frequencies} fields specify $(k,\ell)$, $(m,n)$, and the four ordered rows. For a signed system, the tuple
\[
 (i,j,k,l,\varepsilon,\eta)
\]
means precisely the two equations in~\eqref{eq:roles}. The tuple, rather than an unqualified integer system number, is the portable identifier: the two-parameter and D1 catalogs use different sign-enumeration orders.

A sparse bivariate term is stored as \texttt{[[a,b],c]} and means $ct^au^b$; a univariate resultant term \texttt{[[a],c]} means $ct^a$. In contrast, \texttt{coefficients\_high\_to\_low} in the factor catalog lists coefficients in descending order, including zeros. The fields \texttt{G1}, \texttt{G2}, \texttt{removed\_t\_power}, \texttt{F}, and \texttt{G} are checked through
\[
 G_1=t^jF,\qquad G_2=t^jG.
\]
The branch $t=0$ is checked separately whenever $j>0$. A nonzero integer \texttt{constant} and the complete factor list must multiply to the recorded resultant, with every multiplicity retained. A no-root record supplies both its witness prime and the polynomial's value at every residue.

The acceptance check does not assume factor irreducibility and does not infer absence of rational solutions from a bounded search. It proves the resultant identity using the independent degree bound of Lemma~\ref{lem:elimination}, then checks the modular no-root witnesses and the remaining fibers.

\subsection{A complete worked elimination}
Consider $(k,\ell;m,n)=(2,1;1,0)$ and the assignment
\[
 (i,j,k,l,\varepsilon,\eta)=(0,1,2,3,-1,-1).
\]
With the numerators in~\eqref{eq:twoinputs}, this gives $F=E_0+E_2+E_3$ and $G=E_1+E_2-E_3$. No common $t$ is removed. As quadratics in $u$,
\begin{align*}
 F&=-2t(t^2+1)u^2+(4-12t^2)u+2t(5-3t^2),\\
 G&=2t(3t^2-5)u^2+4t^2(3-t^2)u+2t(t^2+1).
\end{align*}
Writing their coefficients as $f_2,f_1,f_0$ and $g_2,g_1,g_0$, the certificate identity is
\begin{equation}\label{eq:workedresultant}
 \det\begin{pmatrix}
 f_2&f_1&f_0&0\\0&f_2&f_1&f_0\\
 g_2&g_1&g_0&0\\0&g_2&g_1&g_0
 \end{pmatrix}=64t^2J(t),
\end{equation}
where
\[
 J(t)=3t^{12}+26t^{10}-251t^8+716t^6-803t^4+250t^2-5.
\]
Modulo $7$, its values at $0,1,\ldots,6$ are
\[
 (2,6,3,5,5,3,6).
\]
Its leading coefficient is $3\not\equiv0\pmod7$. Thus $J$ has no rational root by the denominator argument in Lemma~\ref{lem:elimination}. The only possible finite rational $t$ in this example is zero, where two offset magnitudes coincide. Both parameter-infinity boundaries are excluded by~\eqref{eq:collision}. The same acceptance steps apply to every record, with the factor lists and rational fibers shown in Table~\ref{tab:representatives}.

\section{The squarefree three-prime configuration list}\label{app:orbits}
The classification in Proposition~\ref{prop:orbits} concerns four-element sets of nonzero projective frequency vectors, not solutions of the circle equations. Choose the representative of each pair $\{v,-v\}$ whose first nonzero coordinate is positive. Of the $\binom{13}{4}=715$ four-subsets, retain exactly those for which every coordinate is nonzero in at least three rows. This gives 177 subsets.

The acting group is the group of all signed coordinate permutations, of order $2^3\cdot3!=48$. Its action on projective vectors has kernel $\{I,-I\}$. The stabilizers below are computed in the full group of order 48, not its effective quotient of order 24; consequently each row satisfies
\[
 |\text{orbit}|\,|\text{stabilizer}|=48.
\]
Rows within a representative are unordered for classification, but their displayed order is used in every explicit matrix identity. The equalities $E3=E4$ and $F{-}1=F{-}3$ identify historical duplicate labels, not additional orbits.

{\small\setlength{\tabcolsep}{4pt}\renewcommand{\arraystretch}{1.18}
\begin{longtable}{lccccrr}
\caption{All admissible four-sets under signed coordinate permutations. The vector order is fixed for matrix identities. Stabilizers refer to the group of order 48.}\label{tab:orbits}\\
\toprule
Label & $v_0$ & $v_1$ & $v_2$ & $v_3$ & Size & Stab.\\
\midrule\endfirsthead
\toprule
Label & $v_0$ & $v_1$ & $v_2$ & $v_3$ & Size & Stab.\\
\midrule\endhead
$A$ & $(1,1,1)$ & $(1,1,-1)$ & $(1,-1,1)$ & $(1,-1,-1)$ & 1 & 48\\
$B1$ & $(1,1,-1)$ & $(1,-1,1)$ & $(1,-1,-1)$ & $(1,1,0)$ & 12 & 4\\
$B2$ & $(1,1,1)$ & $(1,1,-1)$ & $(1,-1,-1)$ & $(1,1,0)$ & 12 & 4\\
$C$ & $(1,1,1)$ & $(1,1,-1)$ & $(1,-1,1)$ & $(0,0,1)$ & 12 & 4\\
$D1$ & $(1,1,1)$ & $(1,1,-1)$ & $(1,1,0)$ & $(1,0,1)$ & 24 & 2\\
$D2$ & $(1,1,1)$ & $(1,-1,-1)$ & $(1,1,0)$ & $(1,0,1)$ & 12 & 4\\
$D3$ & $(1,1,-1)$ & $(1,-1,1)$ & $(1,1,0)$ & $(1,0,1)$ & 12 & 4\\
$D4$ & $(1,1,-1)$ & $(1,-1,-1)$ & $(1,1,0)$ & $(1,0,1)$ & 24 & 2\\
$E1$ & $(1,1,1)$ & $(1,1,-1)$ & $(1,1,0)$ & $(0,0,1)$ & 6 & 8\\
$E2$ & $(1,-1,1)$ & $(1,-1,-1)$ & $(1,1,0)$ & $(0,0,1)$ & 6 & 8\\
$E3=E4$ & $(1,1,1)$ & $(1,-1,1)$ & $(1,1,0)$ & $(0,0,1)$ & 24 & 2\\
$F{-}1=F{-}3$ & $(1,1,1)$ & $(1,1,0)$ & $(1,0,1)$ & $(0,1,-1)$ & 12 & 4\\
$F{-}2$ & $(1,-1,-1)$ & $(1,1,0)$ & $(1,0,1)$ & $(0,1,-1)$ & 4 & 12\\
$F{+}1$ & $(1,1,1)$ & $(1,1,0)$ & $(1,0,1)$ & $(0,1,1)$ & 4 & 12\\
$F{+}2$ & $(1,1,-1)$ & $(1,1,0)$ & $(1,0,1)$ & $(0,1,1)$ & 12 & 4\\
\midrule
\multicolumn{5}{l}{Total admissible four-sets} & 177 & \\
\bottomrule
\end{longtable}
}


The fifteen listed orbits are pairwise disjoint and exhaust the 177 retained subsets. In particular, their sizes sum to 177. The torus isomorphism between $B2$ and $C$ in Section~\ref{sec:three} uses a larger class of integer coordinate changes and therefore does not contradict their being distinct rows of this finite-group classification. This table makes no claim that all fifteen configurations are excluded.

\section{Local descent witnesses and complete elliptic groups}\label{app:descent}
\subsection{Square classes and primitive covering equations}
For $E:y^2=x^3+Ax^2+Bx$, the descent map at a point with $x\ne0$ is its $x$-coordinate modulo $\Q^{\times2}$. Together with the definitions at $\mathcal O$ and $(0,0)$, its image is contained in the signed squarefree divisors $d$ of $B$~\cite{SilvermanTate,SageDescent}. If $x=d(U/V)^2$, with coprime integers $U,V$ and $V\ne0$, substitution gives
\[
 N^2=dU^4+A U^2V^2+(B/d)V^4,\qquad
 y=\frac{dUN}{V^3}.
\]
The rational number $N$ is an integer, since its square is an integer. The cases $V=0$ and $U=0$ recover the identity class and the class of $B$, respectively. In the local obstruction tests, ``primitive'' means $\gcd(U,V)=1$, so at least one of $U,V$ is a unit at every prime.

The complete candidate and retained square-class data are:
\begin{center}\small\setlength{\tabcolsep}{5pt}
\begin{tabular}{rrlll}
\toprule
$A$ &$B$ & Classes not excluded & Excluded classes & Reason\\
\midrule
$6$ &$1$ &$-1,1$ & none & ---\\
$-12$ &$32$ &$1,2$ &$-2,-1$ & real sign\\
$-3$ &$4$ &$1,2$ &$-2,-1$ & real sign\\
$6$ &$-7$ &$-7,1$ &$-1,7$ & modulo $16$\\
$18$ &$1$ &$-1,1$ & none & ---\\
$-36$ &$320$ &$1,5$ &$-10,-5,-2,-1$ & real sign\\
&& &$2,10$ & modulo $64$\\
\bottomrule
\end{tabular}
\end{center}
Every real-sign row has negative coefficients of $U^4$, $U^2V^2$, and $V^4$, so its right side is strictly negative for a nonzero real pair. A class left in the table is merely not excluded at this step. The image upper bounds, two on each side, already suffice in~\eqref{eq:descentrank} to prove rank zero; no local-to-global converse is assumed.

\subsection{Exhaustive local charts}
For modulus $p^e$, if $V$ is a unit divide the covering equation by $V^4$ and test $V=1$ with all residues $U/V$. Otherwise $U$ is a unit; divide by $U^4$ and test $U=1$ with $V/U$ divisible by $p$. In each case the corresponding rescaling of $N$ is by the square of a unit. Thus the property of being a square is preserved. These two charts exhaust all primitive pairs.

The only finite-modulus obstructions needed here are:
\begin{center}\small
\begin{tabular}{crcll}
\toprule
$(A,B)$ & $d$ & Modulus & $V=1$, all $U$ & $U=1$, even $V$\\
\midrule
$(6,-7)$ & $-1$ &16& $\{7,12,15\}$ & $\{7,15\}$\\
$(6,-7)$ & $7$ &16& $\{7,12,15\}$ & $\{7,15\}$\\
$(-36,320)$ & $2$ &64& $\{32,48,62\}$ & $\{2,50\}$\\
$(-36,320)$ & $10$ &64& $\{6,32,48\}$ & $\{10,58\}$\\
\bottomrule
\end{tabular}
\end{center}
The full square-residue sets are
\begin{align*}
 \operatorname{Sq}(16)&=\{0,1,4,9\},\\
 \operatorname{Sq}(64)&=\{0,1,4,9,16,17,25,33,36,41,49,57\}.
\end{align*}
Every displayed covering-residue set is disjoint from the corresponding square-residue set. The companion checks also enumerate all primitive pairs directly, providing a second check on the normalized-chart calculations.

\subsection{Torsion bounds prime by prime}
Rank zero alone is not a complete point list. For the three original models, the following good-prime reductions bound the rational torsion:
\begin{center}\small
\begin{tabular}{ccl}
\toprule
$(A,B)$ & $(p,\#E(\F_p))$ & Consequence\\
\midrule
$(6,1)$ & $(3,4),(5,8)$ & torsion order divides $4$\\
$(-3,4)$ & $(3,4),(5,4)$ & torsion order divides $4$\\
$(18,1)$ & $(3,4),(7,4)$ & torsion order divides $4$\\
\bottomrule
\end{tabular}
\end{center}
All indicated primes are good because they do not divide $16B^2(A^2-4B)$. For a torsion prime $\lambda\ne p$, the $\lambda$-primary rational torsion injects into $E(\F_p)$. The reduction at $p=3$ bounds the $2$-primary order by four and excludes every other prime except possibly $3$. The second reduction, whose order is not divisible by $3$, excludes $3$-primary torsion as well. This supplies the order-four bound without assuming injectivity of the full torsion group at a prime dividing its order. The four points in Table~\ref{tab:groups} satisfy the equations and are distinct, so they exhaust the rank-zero groups.

\section{D1 normalization records and exceptional points}\label{app:d1checks}
\subsection{Indexing the complete normalization catalog}
For reproducibility, the D1 catalog orders the pairs $i<j$ lexicographically, then the two orders of the complementary pair $(k,l)$, then $(\varepsilon,\eta)$ in
\[
 (1,1),(1,-1),(-1,1),(-1,-1).
\]
Record numbers are zero-based. This convention produces the following partition of all 48 records:
\begin{center}\small
\begin{tabular}{lrp{0.47\linewidth}}
\toprule
$N(r)$ & Count & Zero-based record indices\\
\midrule
$1$ & 16 & 0, 1, 2, 3, 18, 19, 21, 23, 34, 35, 37, 39, 40, 43, 44, 47\\
$r^{4} + 6 r^{2} + 1$ & 4 & 4, 5, 6, 7\\
$r^{4} - 3 r^{2} + 4$ & 4 & 10, 11, 26, 27\\
$4 r^{4} - 3 r^{2} + 1$ & 4 & 13, 15, 29, 31\\
$r^{4} + 18 r^{2} + 1$ & 12 & 16, 17, 20, 22, 32, 33, 36, 38, 41, 42, 45, 46\\
$r^{2} + 9$ & 4 & 8, 9, 24, 25\\
$9 r^{2} + 1$ & 4 & 12, 14, 28, 30\\
\bottomrule
\end{tabular}
\end{center}

The full kernel vectors and common divisors are retained in \path{verification/d1_norm/d1_certificate.json}. The eight residual records are $8,9,12,14,24,25,28,30$; the other forty are excluded as in Corollary~\ref{cor:forty}. In particular, sixteen records are degenerate before any elliptic argument, and twenty-four are excluded using the three elliptic descents.

For record 8, the relation rows are
\[
 R=\begin{pmatrix}1&-1&0&-1\\0&-1&1&1\end{pmatrix}.
\]
Their images under $H(r)$ are
\[
 RH(r)=\begin{pmatrix}0&4r&-q\\2r^2&2r&q\end{pmatrix},
 \qquad q=1+r^2.
\]
Their cross product is $2r(3q,-rq,-4r^2)$, yielding the two norms
\[
 (3q)^2+(-rq)^2=q^2(r^2+9),\qquad
 q^2(r^2+9)-16r^4=r^6-5r^4+19r^2+9.
\]
This shows explicitly why the first square condition and the second discriminant condition must both be retained.

For every record, the primitive kernel entries have gcd one in $\Q[r]$, so they do not vanish simultaneously at any parameter. The factors $r^2-3$ and $3r^2-1$ in the removed scalar have no rational roots: modulo $5$ their value lists are $(2,3,1,1,3)$ and $(4,2,1,1,2)$, with nonzero leading coefficients. The only rational rank-drop parameter is therefore $r=0$. That entire fiber, and $r=\infty$, have repeated first two sine magnitudes and are excluded directly.

\subsection{A short squarefreeness witness and the infinity chart}
The genus assertion in Lemma~\ref{lem:pullback} can be checked without a long integer Euclidean remainder sequence. Modulo $5$,
\begin{equation}\label{eq:Pmodfive}
 P(h)\equiv(h^2-2)(h^2+2)(h^4-2)(h^4+2).
\end{equation}
These factors are squarefree and pairwise coprime over $\F_5$. For instance, a common root of a quadratic and a quartic would have $h^4=4$ and also $h^4=2$ or $-2$, a contradiction in $\F_5$. The derivatives exclude repeated roots within each factor. Since $P$ is monic, this proves squarefreeness over $\Q$. The degree-twelve double cover has twelve simple geometric branch points and hence genus five.

At infinity use $t=1/h$ and $\eta=Y/h^6$. Reciprocity of $P$ gives the local equation $\eta^2=P(t)$. At $t=0$ the two points are $(0,1)$ and $(0,-1)$, both smooth since $2\eta\ne0$. This verifies the two rational points at infinity and distinguishes them from the single coordinate value $h=\infty$. The eight points in $\mathcal K$ are exactly the fibers over $h=0,1,-1,\infty$; this says nothing about the existence of points over other rational parameters.

\subsection{All reconstruction collisions}
Let $e_0,e_1,e_2,e_3$ be the four polynomials in parentheses in~\eqref{eq:reconstructedsines}. Their six pairwise differences and six pairwise sums are:
\begin{center}\small
\begin{tabular}{cll}
\toprule
$(i,j)$ & $e_i-e_j$ & $e_i+e_j$\\
\midrule
$(0,1)$ & $- 48 h^{2}$ & $2 \left(h^{4} - 38 h^{2} + 1\right)$\\
$(0,2)$ & $- 96 h^{2}$ & $2 \left(h^{4} - 14 h^{2} + 1\right)$\\
$(0,3)$ & $h^{4} - 14 h^{2} + 1$ & $h^{4} - 110 h^{2} + 1$\\
$(1,2)$ & $- 48 h^{2}$ & $2 \left(h^{4} + 10 h^{2} + 1\right)$\\
$(1,3)$ & $h^{4} + 34 h^{2} + 1$ & $h^{4} - 62 h^{2} + 1$\\
$(2,3)$ & $h^{4} + 82 h^{2} + 1$ & $h^{4} - 14 h^{2} + 1$\\
\bottomrule
\end{tabular}
\end{center}

The entries themselves are $h^4-62h^2+1$, $h^4-14h^2+1$, $h^4+34h^2+1$, and $-48h^2$. After extracting nonzero constants and powers of $h$, every nonconstant expression in this list or the table is $h^4+ch^2+1$ with
\[
 c\in\{-110,-62,-38,-14,10,34,82\}.
\]
The rational-root theorem restricts a rational zero to $h=\pm1$, where the value is $c+2\ne0$. For $h\notin\{0,\pm1\}$, the common multiplier $L=(1-h^2)/D(h)$ is nonzero and $D(h)>0$. Thus all four signed offsets are nonzero and their magnitudes are pairwise distinct.

This last check is needed for the converse in Theorem~\ref{thm:main-d1}; solving the two linear relations alone would not suffice. Together with the conic inverse and the infinity chart, it accounts for every rational exceptional parameter used in that theorem. It does not prove Conjecture~\ref{conj:points}.


\begin{thebibliography}{99}
\small
\bibitem{AuelSinger} A. Auel and B. Singer, \emph{The algebraic geometry of 3-by-3 magic squares of squares}, arXiv:2609.09351v1 (2026). \url{https://arxiv.org/abs/2609.09351v1}.
\bibitem{Bremner} A. Bremner, On squares of squares, \emph{Acta Arith.} \textbf{88} (1999), no.~3, 289--297. \href{https://doi.org/10.4064/aa-88-3-289-297}{doi:10.4064/aa-88-3-289-297}.
\bibitem{Brown511} K. Brown, \emph{Discordance Impedes Square Magic}, MathPages, undated. \url{https://www.mathpages.com/home/kmath511/kmath511.htm}. Accessed 2 October 2026.
\bibitem{Brown417} K. Brown, \emph{Magic Square of Squares}, MathPages, Proposition~1, undated. \url{https://www.mathpages.com/home/kmath417/kmath417.htm}. Accessed 2 October 2026.
\bibitem{BTV} N. Bruin, J. Thomas, and A. V\'arilly-Alvarado, Explicit computation of symmetric differentials and its application to quasihyperbolicity, \emph{Algebra Number Theory} \textbf{16} (2022), no.~6, 1377--1405. \href{https://doi.org/10.2140/ant.2022.16.1377}{doi:10.2140/ant.2022.16.1377}.
\bibitem{Cain} O. Cain, \emph{Gaussian Integers, Rings, Finite Fields, and the Magic Square of Squares}, arXiv:1908.03236v2 (2019). \url{https://arxiv.org/abs/1908.03236v2}.
\bibitem{Houston} R. Houston, \emph{Squares of squares, and the group of rational points on the circle}, Bosker Blog, 2 August 2016, formulations P4 and W1. \url{https://bosker.wordpress.com/2016/08/02/squares-of-squares-and-the-group-of-rational-points-on-the-circle/}. Accessed 2 October 2026.
\bibitem{MSSArchive} Magic Squares of Squares project, \emph{Refined prime-power block balance}, public unrefereed research note, Sections~3--4; project maintained by U. Schwarz. Archive commit dated 20 August 2026. \url{https://github.com/uwe-schwarz/magic-squares-of-squares/blob/220c912a01186b406128abb6b6e74fc4d9cc4f2f/research/prime-support/block-balance.md}.
\bibitem{MSSTwoBlock} Magic Squares of Squares project, \emph{Excluding every two-block center}, public unrefereed research note (2026). \url{https://github.com/uwe-schwarz/magic-squares-of-squares/blob/main/research/prime-support/two-block-exponent-one.md}. Accessed 2 October 2026.
\bibitem{SymPy} A. Meurer et al., SymPy: symbolic computing in Python, \emph{PeerJ Comput. Sci.} \textbf{3} (2017), e103. \href{https://doi.org/10.7717/peerj-cs.103}{doi:10.7717/peerj-cs.103}. The computations here use version 1.14.0.
\bibitem{BHPV} W. P. Barth, K. Hulek, C. A. M. Peters, and A. Van de Ven, \emph{Compact Complex Surfaces}, 2nd ed., Ergebnisse der Mathematik (3) \textbf{4}, Springer, Berlin, 2004.
\bibitem{Bremner2001} A. Bremner, On squares of squares II, \emph{Acta Arith.} \textbf{99} (2001), no.~3, 289--308.
\bibitem{Lang} S. Lang, Hyperbolic and Diophantine analysis, \emph{Bull. Amer. Math. Soc.} \textbf{14} (1986), no.~2, 159--205.
\bibitem{Pardini} R. Pardini, Abelian covers of algebraic varieties, \emph{J. Reine Angew. Math.} \textbf{417} (1991), 191--213.
\bibitem{Robertson} J. P. Robertson, Magic squares of squares, \emph{Math. Mag.} \textbf{69} (1996), no.~4, 289--293. \href{https://doi.org/10.1080/0025570X.1996.11996457}{doi:10.1080/0025570X.1996.11996457}.
\bibitem{SageDescent} SageMath developers, \emph{Descent on elliptic curves over $\Q$ with a 2-isogeny}, SageMath reference manual. \url{https://doc.sagemath.org/html/en/reference/arithmetic_curves/sage/schemes/elliptic_curves/descent_two_isogeny.html}. Accessed 2 October 2026.
\bibitem{SilvermanTate} J. H. Silverman and J. T. Tate, \emph{Rational Points on Elliptic Curves}, 2nd ed., Undergraduate Texts in Mathematics, Springer, Cham, 2015. \href{https://doi.org/10.1007/978-3-319-18588-0}{doi:10.1007/978-3-319-18588-0}.
\bibitem{Weisenberg} D. Weisenberg, Some Thoughts on The $3\times3$ Magic Square of Squares Problem, \emph{Rose-Hulman Undergrad. Math. J.} \textbf{24} (2023), no.~1, Article~7. \url{https://scholar.rose-hulman.edu/rhumj/vol24/iss1/7/}.
\bibitem{Wolird} C. Wolird, \emph{A New Transformation of the Magic Square of Squares}, arXiv:2310.12164v1 (2023). \url{https://arxiv.org/abs/2310.12164v1}.
\end{thebibliography}

\end{document}
